
\documentstyle[amsfonts,12pt]{article}
%%%%%%%%%%%%%%%%%%%%%%%%%%%%
%TCIDATA{TCIstyle=Article/ART4.LAT,jart,sw20jart}

%\input tcilatex
%\QQQ{Language}{
%American English
%}

\begin{document}

\author{David Joyner}
\title{A hitch-hiker's guide to invariant distributions on SL(2) over a p-adic field}
\date{11-2-94}
\maketitle

\begin{abstract}
In this expository article, we write down the admissible, unitary, and
tempered duals on $SL(2)$, using the explicit parameterization of
Gelfand-Graev (and results of Sally-Shalika). This allows us to determine
the Tadi\'c decompositions and natural Borel structures on the dual spaces.
Then we classify the image of the Schwartz space (resp., the Hecke algebra)
under the Fourier transform and define the Fourier transform of any tempered
(resp., admissible) distribution on $SL(2)$. The Fourier transform of a
tempered distribution can be made considerable more explicit than that of an
admissible distribution. In any case, these results are then applied to an
understanding of supertempered distributions on $SL(2)$.
\end{abstract}

\tableofcontents


\section{Introduction}

This paper is a essentially an expository article. It is not by any means a
complete survey nor a guide to the literature. It is written to work out the
ideas in [J] in a special case. Proofs and pointers to the literature are
often omitted, although everything should be in one of the references listed
at the end of the paper. In order to make this paper appeal to a broader
audience (and to make it a more useful reference for myself and hopefully
others), lots of definitions are included and I've tried my best to (a) be
precise, (b) be as non-technical as possible, and (c) preserve the intrinsic
beauty of the subject (which is largely due to Harish-Chandra). To be
honest, some of the definitions and examples in \S \S 7-9 seem to be
unpublished, so this article is not entirely expository. I hope what little
new there is in those sections does not remind anyone of the old joke about
the PhD thesis (it was both interesting and new, but the interesting parts
where not new and the original parts were not interesting).
%Questions,
%comments or constructive flames may be sent to wdj@usna.navy.mil

\subsection{Basic notation}

Let $G=SL(2,F)$, where $F$ is a p-adic field with uniformizer $\pi =\pi _F$,
ring of integers $O_F,$ residual characteristic $p=char(O_F/\pi O_F),$ $%
q=\#(O_F/\pi O_F),$ and normalized valuation $|...|=|...|_F.$ We always
assume $p>2.$ Let $\varepsilon $ denote a fixed non-square unit so that $%
1,\varepsilon ,\pi ,\varepsilon \pi $ form a complete set of representatives
of $F^{\times }/(F^{\times })^2.$ The subgroups 
\[
U_0=O_F^{\times },\;U_n=1+\pi ^nO_F,\;\;\;\;n>0,\; 
\]
form a basis of open neighborhoods of the identity in $F^{\times }$,
generating the topology on $F^{\times }$. The subgroups 
\[
K_0=SL(2,O_F),\;K_n=\left\{ \left[ 
\begin{array}{cc}
a & b \\ 
c & d
\end{array}
\right] \;\;|\;\;a,d\in 1+\pi ^nO_F,\;b,c\in \pi ^nO_F\right\}
,\;\;\;\;n>0,\; 
\]
form a basis of open neighborhoods of the identity in $G$, generating the
topology on $G$. A {\bf character} or {\bf representation} of $F^{\times }$
or of $G$ will always mean a continuous multiplicative (not necessarily
unitary) homomorphism into the automorphisms of a complex vector space,
where the domain gets the ``compact-open'' topology above and the range gets
the discrete topology.

Occasionally we use {\bf Vinogradov's notation}: given two functions $f(x)$,
complex-valued and perhaps depending on some parameter $\lambda $, and $%
g(x)>0$, we write 
\[
f(x)<<_\lambda g(x), 
\]
if there is a constant $c>0,$ which may implicitly depend on $\lambda $,
such that $|f(x)|\leq cg(x)$ for all $x$ in some range which depends on the
context.

\section{Background on the unitary dual and the admissible dual}

\subsection{Definitions}

Let $V$ denote a complex vector space. A representation $\pi :G\rightarrow
Aut(V)$ is called {\bf admissible} if

(1) every $v\in V$ belongs to $V^K=stab_V(K)$, for some compact open $%
K\subset \subset G,$ (i.e., $\pi $ is {\bf smooth}),

(2) for every compact open $K\subset \subset G,$ $\dim (V^K)<\infty .$

Let $V^{\prime }$ denote the dual space of all complex-valued linear
functionals on $V$ and let $<v,v^{\prime }>=v^{\prime }(v)$ denote the
evaluation pairing on $V\times V^{\prime }$. Let $V_\infty ^{\prime }$
denote the subspace of $V^{\prime }$ of vectors fixed by some compact open
subgroup. Define the {\bf contragredient} representation of a smooth $\pi
:G\rightarrow End(V)$ to be the representation $\pi ^{\prime }:G\rightarrow
Aut(V_\infty ^{\prime })$ satisfying $<v,v^{\prime }.=<\pi (g)v,\pi ^{\prime
}(g)v^{\prime }>$, for all $v\in V,v^{\prime }\in V_\infty ^{\prime }$ , and 
$g\in G.$ It is known that if a representation $\pi :G\rightarrow Aut(V)$ is
admissible then its contragredient representation $\pi ^{\prime
}:G\rightarrow Aut(V_\infty ^{\prime })$ is admissible. A {\bf matrix
coefficient} of a representation $\pi :G\rightarrow Aut(V)$ is a function on 
$G$ of the form $g\longmapsto <\pi (g)v,v^{\prime }>,$ for some $v\in
V,v^{\prime }\in V_\infty ^{\prime }.$ The {\bf space of matrix coefficients}
of $\pi $ is denoted by ${\cal A}(\pi ).$

An admissible representation $\pi :G\rightarrow Aut(V)$ is called {\bf %
square-integrable} or {\bf in the discrete series} if ${\cal A}(\pi )\subset
L^2(G).$ It is called {\bf cuspidal} if every function $f\in {\cal A}(\pi )$
has compact support modulo the center and satisfies, for each nilpotent
radical $N$ of a proper parabolic subgroup of $G$, $\int_Nf(xn)dn=0$, for
all $x\in G$. Clearly, a cuspidal representation is square-integrable. A
representation $\pi :G\rightarrow End(V)$ is called {\bf unitarizable} if $V$
has a positive-definite $G$-invariant inner product $(*,*)$ (that is, $(\pi
(g)v,\pi (g)v^{\prime })=(v,v^{\prime })$ for all $v,v^{\prime }\in V$). The
completion of $V$ with respect to this inner product is of course a $G$%
-invariant Hilbert space $H$. Any representation $\pi :G\rightarrow Aut(H),$
for some Hilbert space $H$ whose inner product is $G$-invariant, is called 
{\bf unitary. }It is known that every square-integrable representation is
unitarizable.

The following theorem was proven by Harish-Chandra subject one of two
statements, one of which was later proven by Bernstein, and the other was
later proven by Clozel.

{\bf Theorem} (Harish-Chandra, Bernstein, Clozel) {\em Any unitary
representation }$\pi :G\rightarrow Aut(H)${\em , has the property that there
is a dense }$G${\em -invariant subspace }$V=H_\infty ${\em \ such that }$%
(\pi ,V)${\em \ is admissible.}

We say that a representation $\pi _1:G\rightarrow Aut(V_1)$ is {\bf %
equivalent} to a representation $\pi _2:G\rightarrow Aut(V_2)$ if there is a
non-zero linear transformation $A:V_1\rightarrow V_2$ such that $A\pi
_1(g)v=\pi _2(g)Av$ for all $g\in G,$ $v\in V_1.$ The map $A$ is called an 
{\bf intertwining map}. The {\bf unitary dual} of $G$ is the set $\hat G_u$
of all equivalence classes of irreducible unitary representations of $G$ and
the {\bf admissible dual} of $G$ is the set $\hat G_{ad}$ of all equivalence
classes of irreducible admissible representations of $G.$ The only finite
dimensional irreducible admissible representation of $G$ is the trivial
representation. As a corollary to the above theorem, we have a natural
inclusion $\hat G_u\subset \hat G_{ad}.$

Each admissible representation $\pi :G\rightarrow Aut(V)$ gives rise to a
representation representation $\pi :C_c^\infty (G)\rightarrow End(V)$ of the
Hecke algebra $C_c^\infty (G)$ of all locally constant functions with
compact support on $G$: 
\[
\pi (f)v=\int_Gf(g)\pi (g)v\;dg,\;\;\;v\in V,\;\;f\in C_c^\infty (G), 
\]
where $dg$ denotes a Haar measure on $G$. The linear transformation $\pi (f)$
is sometimes called the {\bf operator-valued Fourier transform} of $f$ at $%
\pi .$ For admissible $\pi $ it is known that $\pi (f)$ is finite rank as an
element of $End(V)$. Let $||\pi (f)||$ denote its operator norm and let 
\[
||f||_{*}=\sup_{\pi \in \hat G_u}||\pi (f)||. 
\]
The completion of $C_c^\infty (G)$ with respect to this norm $||...||_{*}$
is a $C^{*}-$ algebra denoted by $C^{*}(G)$.

Let $A$ denote a $C^{*}$-algebra. A primitive ideal of $A$ is the kernel of
an irreducible representation of $A$. Put the Jacobson topology on the set $%
\Pr (A)$ of all primitive ideals of $A$. Let $A^{\wedge }$ denote the set of
equivalence classes of irreducible representations of $A$ on the bounded
operators on a Hilbert space. There is a canonical surjection $A^{\wedge
}\hookrightarrow \Pr (A)$ induced by $\pi \longmapsto \ker (\pi ).$ Give $%
A^{\wedge }$ the smallest topology such that this surjection is continuous.
Call this the {\bf Jacobson} topology. This topology is $T_0$.

We define a topology on $\hat G_u$ using the following

{\bf Theorem} (Dixmier): {\em There is a canonical bijection between the set 
}$\hat G_u${\em \ and the set }$C^{*}(G)^{\wedge }${\em \ of equivalence
classes of irreducible representations of }$C^{*}(G).$

We give $\hat G_u$ the smallest topology such that this bijection is
continuous. Call this the {\bf Jacobson topology} on $\hat G_u.$

A {\bf CCR algebra} is a $C^{*}-$algebra $A$ such that every irreducible
representation $\pi :A\rightarrow End(H)$ is completely continuous (i.e.,
compact). It is known, by a result of Fell, that if $A$ is a CCR algebra and
if we topologize $A^{\wedge }$ using the Jacobson topology as above then the
dual space $A^{\wedge }$ is a $T_1$ topological space.

{\bf Theorem} (Bernstein) $C^{*}(G)${\em \ is a CCR algebra.}

Following Mackey, we say that the dual space $A^{\wedge }$ is {\bf smooth}
if there is a Borel structure on it which is Borel isomorphic to the Borel
structure on the real line.

{\bf Theorem} (Fell, Mackey) {\em If }$A${\em \ is a separable CCR algebra
then }$A^{\wedge }${\em \ has a smooth dual. Moreover, the Borel structure
may be taken to be that generated by the Jacobson topology on }$A^{\wedge }.$

{\sc Question}: Is there a natural topology on $\hat G_{ad}$ which generates
a smooth Borel structure and which, in turn, induces the Borel structure on $%
\hat G_u$ via the inclusion $\hat G_u\subset \hat G_{ad}$?

We shall see that the answer is yes for $SL(2).$

\subsection{The Tadi\'c decomposition}

Let $A$ denote the diagonal subgroup of $G$, and let ${\cal L}(G)$ denote
the set of all standard Levi's (with respect to $A$), up to associates.
Therefore, ${\cal L}(G)=\{A,G\}.$ For $M\in {\cal L}(G),\;$we call a
representation $\sigma :M\rightarrow Aut(V)$ for which $A(\sigma )\subset
C_c^\infty (M)$ a {\bf cuspidal} representation of $M$. For $M\in {\cal L}%
(G),\;$let $M_c^{\wedge }$ denote the set of equivalence classes of cuspidal
representations $\sigma :M\rightarrow Aut(V)$, and let $M_{u,c}^{\wedge }$
denote the set of equivalence classes of cuspidal representations $\sigma
:M\rightarrow Aut(W)$ which has unitary central character $\omega _\sigma $
(i.e., $|\omega _\sigma (z)|=1$ for all $z$ belonging to the center of $M$).
For example, if $M=A$ then 
\[
\begin{array}{c}
A_{u,c}^{\wedge }=\{\chi :A\rightarrow {\bf C}^{\times }\;|\; 
\begin{array}{c}
\chi ( 
\begin{array}{cc}
a & 0 \\ 
0 & a^{-1}
\end{array}
)=\chi ^{\prime }(a),\;where\;\chi ^{\prime }(a)=\chi _0(\frac a{|a|})\cdot
|a|^{i\theta }, \\ 
\chi _0\in (O_F^{\times })^{\wedge },\;-\frac \pi {\log (q)}\leq \theta
<\frac \pi {\log (q)}\}
\end{array}
\\ 
\;\cong (F^{\times })_u^{\wedge }, \\ 
A_c^{\wedge }=\{\chi :A\rightarrow {\bf C}^{\times }\;|\; 
\begin{array}{c}
\chi ( 
\begin{array}{cc}
a & 0 \\ 
0 & a^{-1}
\end{array}
)=\chi ^{\prime }(a),\;where\;\chi ^{\prime }(a)=\chi _0(\frac a{|a|})\cdot
|a|^{t+i\theta }, \\ 
\chi _0\in (O_F^{\times })^{\wedge },\;t\in {\bf R},\;-\frac \pi {\log
(q)}\leq \theta <\frac \pi {\log (q)}\}
\end{array}
\\ 
\cong (F^{\times })^{\wedge },
\end{array}
\]
where, for any commutative topological group $H$, $H^{\wedge }$ denotes the
set of continuous complex-valued (not necessarily unitary) characters of $H$.

Let $X_{un}(M)$ denote the group of all unramified characters of $M$. For
example, if $M=A$ then we may either write $\chi \in X_{un}(A)$ as 
\[
\chi (x)=\chi _0({\frac x{|x|}})|x|^{t+i\theta }, 
\]
with $\chi _0\in (O_F^{\times })^{\wedge },\;t\in {\bf R}$, $\chi |_{1+\pi 
{\cal O}_F}=1$, and $-\frac \pi {\log (q)}\leq \theta <\frac \pi {\log (q)}$%
, {\em or} as 
\[
\chi (x)=z^{\nu _F(x)}, 
\]
where $z\in {\bf C}^{\times }$ and $\nu _F$ denotes the usual valuation map
(i.e., for all $x\in F^{\times }$, $|x|=q^{-\nu _F(x)}$). Therefore, we may
either identify $X_{un}(A)$ as a countable union of cylinders ${\bf C}/{%
\frac{2\pi }{\log (q){\bf Z}}}$ {\em or} as ${\bf C}^{\times }$.

Let $(\sigma ,W)\in M_c^{\wedge }$ and let $I_M(\sigma ):G\rightarrow Aut(V)$
denote the {\bf unitarily induced representation}: the representation of $%
G\; $by right translation on 
\[
V=\{f:G\rightarrow W\;|\; 
\begin{array}{c}
(1)\;\;f(mg)=\delta _M(m)^{1/2}\sigma (m)f(g),\;\; \\ 
\;\forall g\in G,\;m\in M \\ 
(2)\;\;{\rm for\;some\;open\;subgp}\;K\subset \subset G,\;f(gk)=f(g),\;\; \\ 
\forall k\in K,\;g\in G\;
\end{array}
\}. 
\]
Here $\delta _M(m)=|\det (Ad(m)_{{\frak n}})|$ where ${\frak n}$ denotes the
Lie algebra of 
\[
N=\left\{ 
\begin{array}{c}
\{1\},\;\;\;\;\;\;\;\;\;\;\;\;\;\;\;\;\;\;\;\;M=G, \\ 
\{\left[ 
\begin{array}{cc}
1 & x \\ 
0 & 1
\end{array}
\right] \;|\;x\in F\},\;\;\;M=A.
\end{array}
\right. 
\]
so $\delta _A( 
\begin{array}{cc}
a & 0 \\ 
0 & a^{-1}
\end{array}
)=|a|^2.$ Following Tadi\'c, we let 
\[
\begin{array}{c}
\begin{array}{c}
G_{u,\sigma }^{\wedge }=\{{\rm all\;inequivalent\;irreducible\;unitarizable\;%
} \\ 
{\rm constituents\;of\;}I_M(\chi \sigma )\;\;|\;\;\chi \in X_{un}(M)\},
\end{array}
\\ 
\\ 
\begin{array}{c}
G_\sigma ^{\wedge }=\{{\rm all\;inequivalent\;irreducible\;constituents\;}
\\ 
{\rm of\;}I_M(\chi \sigma )\;\;|\;\;\chi \in X_{un}(M)\}.
\end{array}
\end{array}
\]

{\bf Theorem} (Tadi\'c) {\em (a) }$G_u^{\wedge }=\bigcup_{M\in {\cal L}%
(G)}\bigcup_{\sigma \in M_{u,c}^{\wedge }}G_{u,\sigma }^{\wedge },$

{\em (b) Each }$G_{u,\sigma }^{\wedge }${\em \ is both open and closed in }$%
G_u^{\wedge }.$

Each ``component'' $G_{u,\sigma }^{\wedge }$ is not necessarily connected.
We can choose a subset $\tilde{M}$ of $M_{u,c}^{\wedge }$ such that
we have a disjoint union: 
\[
G_u^{\wedge }=\coprod_{M\in {\cal L}(G)}
\coprod_{\sigma \in \tilde{M}}G_{u,\sigma }^{\wedge }, 
\]

We shall see that for the admissible dual there is an analogous
``decomposition'' (this is due to several people, Jacquet in particular [C],
and pre-dates the decomposition above) 
\[
G_{ad}^{\wedge }=\bigcup_{M\in {\cal L}(G)}
\bigcup_{\sigma \in M_c^{\wedge}}G_\sigma ^{\wedge }. 
\]
We will see that we may impose a Borel structure on the admissible dual as
well.

\section{Borel structure on the Tadi\'c components}

Let $V=C(F)_\infty \;$denote the vector space of locally constant
complex-valued functions on $F$which are bi-$K$-invariant for some compact
open subgroup $K\subset \subset F^{\times }.$

\subsection{Principal series, reducible principal series}

For $\chi \in (F^{\times })^{\wedge },$define $T_\chi :G\rightarrow Aut(V)$
by 
\[
(T_\chi (g)f)(x)=f(\frac{dx+b}{cx+a})\chi
(cx+a)|cx+a|^{-1},\;\;\;\;\;g=\left[ 
\begin{array}{cc}
a & b \\ 
c & d
\end{array}
\right] \in G. 
\]

Note that this representation is smooth.

{\bf Lemma} (Gelfand-Graev):{\em \ Let }$\chi ,\chi _1,\chi _2\in (F^{\times
})^{\wedge }$

{\em (1) If }$\chi _1\neq \chi _2^{-1},\chi _2${\em \ then }$T_{\chi _1}$%
{\em \ is not equivalent to }$T_{\chi _2},$

{\em (2) }$T_\chi \cong T_{\chi ^{-1}},$

{\em (3) }$T_\chi ${\em \ is unitary if }$\chi ${\em \ is,}

{\em (4) }$T_\chi ${\em \ is irreducible if and only if }$\chi \neq
sign_\tau ,${\em \ for each }$\tau =\varepsilon ,\pi ,\varepsilon \pi ${\em %
, where }$sign_\tau ${\em \ denotes the non-trivial quadratic character on }$%
F^{\times }${\em \ for which }$sign_\tau (\tau )=1.$

{\em (5) if }$\chi =sign_\tau ,${\em \ for some }$\tau =\varepsilon ,\pi
,\varepsilon \pi ${\em , then }$T_\chi ${\em \ decomposes into a direct sum
of two irreducible inequivalent unitary subrepresentations }$T_\chi
^{+},\;T_\chi ^{-}.$

The irreducible representations in (4) are called {\bf principal series
representations}. The irreducible representations in (5) are called {\bf %
reducible principal series representations}. They, along with the special
representation, are called {\bf limits of discrete series}.

For $\chi \in (F^{\times })_u^{\wedge }\;$ramified, let 
\[
\begin{array}{c}
G_{u,\chi }^{\wedge }=\{ 
\begin{array}{c}
T_{\chi ^{\prime }}\;,\; {\rm \;if\;}\chi ^{\prime }\neq sign_\tau ,\;{\rm %
for\;each}\;\tau =\pi ,\varepsilon \pi \\ 
\{T_{\chi ^{\prime }}^{+}\;,\;T_{\chi ^{\prime }}^{-}\},\;\;{\rm \;if\;}\chi
^{\prime }=sign_\tau ,\;{\rm for\;some}\;\tau =\pi ,\varepsilon \pi
\end{array}
\\ 
\;|\;\;\;\chi ^{\prime }(x)=\chi (x)z^{\nu _F(x)},\;\;|z|=1,\;\;z\in {\bf C}%
^{\times }/\tilde\},
\end{array}
\]
where $z_1\sim z_2$ if and only if \ $\chi (x)z_1^{\nu _F(x)}=\chi
(x)^{-1}z_2^{-\nu _F(x)}$ , $\forall x\in F^{\times }.$ This is either
isomorphic to a half-circle union two doubleton sets (if the ramified
characters of order two occur in the orbit of $\chi $ under $X_{un}(A)$) or
to a circle (if the ramified characters of order two don't occur in the
orbit of $\chi $ under $X_{un}(A)$), where $\{\pm 1\}$ acts via the
inversion map. All the $T_{\chi ^{\prime }}$ occurring in the above set are
irreducible and inequivalent For example, if $\chi =sign_\pi $ (resp., $\chi
=sign_{\varepsilon \pi }$) then $\chi ^{\prime }=sign_{\pi \varepsilon }$
(resp., $\chi ^{\prime }=sign_\pi $) is also in it's orbit under the action
of the unramified characters of $A.$

Note that the set $G_{u,\chi }^{\wedge }$ and the set $G_{u,\chi
^{-1}}^{\wedge }$ contain equivalent representations. We write $(F^{\times
})_u^{\wedge }/\sim$ for the set of unitary characters modulo the
equivalence relation: $\chi \sim \chi ^{\prime }$ if and only if $\chi
^{\prime }=\chi ^{-1}$.

For $\chi \in (F^{\times })^{\wedge }\;$ramified, let 
\[
\begin{array}{c}
G_{ad,\chi }^{\wedge }=\{ 
\begin{array}{c}
T_{\chi ^{\prime }}\;,\; {\rm \;if\;}\chi ^{\prime }\neq sign_\tau ,\;{\rm %
for\;each}\;\tau =\pi ,\varepsilon \pi \\ 
\{T_{\chi ^{\prime }}^{+}\;,\;T_{\chi ^{\prime }}^{-}\},\;\;{\rm \;if\;}\chi
^{\prime }=sign_\tau ,\;{\rm for\;some}\;\tau =\pi ,\varepsilon \pi
\end{array}
\\ 
\;|\;\;\;\chi ^{\prime }(x)=\chi (x)z^{\nu _F(x)}\;\;z\in {\bf C}^{\times }/%
\tilde\},
\end{array}
\]
where the equivalence action is as above. We can also write this as

\[
\begin{array}{c}
G_{ad,\chi }^{\wedge }=\{ 
\begin{array}{c}
T_{\chi ^{\prime }}\;,\; {\rm \;if\;}\chi ^{\prime }\neq sign_\tau ,\;{\rm %
for\;each}\;\tau =\pi ,\varepsilon \pi \\ 
\{T_{\chi ^{\prime }}^{+}\;,\;T_{\chi ^{\prime }}^{-}\},\;\;{\rm \;if\;}\chi
^{\prime }=sign_\tau ,\;{\rm for\;some}\;\tau =\pi ,\varepsilon \pi
\end{array}
\\ 
\;|\;\;\;\chi ^{\prime }(x)=\chi (x)\chi _0(\frac x{|x|})|x|^{t+i\theta
},\;\; {\rm some} \\ 
\;\;\;\;\;\;\;\;\;\;\;\;(\chi _0,t,\theta )\in \{(O_F^{\times
})_{un}^{\wedge }\times {\bf R}\times [-\frac \pi {\log (q)},\frac \pi {\log
(q)})\}/\tilde\},
\end{array}
\]
where $(\chi _0,t,\theta )\sim (\chi _0^{\prime },t^{\prime },\theta
^{\prime })$ if and only if \ $\chi _0^{\prime }=\chi _0^{-1},\;t^{\prime
}=-t,$ and $\theta ^{\prime }=-\theta .$ All the $T_{\chi ^{\prime }}$
occurring in the above set are irreducible and inequivalent. This is either
isomorphic to ${\bf C}^{\times }/\{\pm 1\}$ with two double points (if the
ramified characters of order two occur in the orbit of $\chi $ under $%
X_{un}(A)$) or to $\{{\bf C}^{\times }\}/\{\pm 1\}$ (if the ramified
characters of order two don't occur in the orbit of $\chi $ under $X_{un}(A)$%
), where $\{\pm 1\}$ acts via the inversion map.

Note that the set $G_{ad,\chi }^{\wedge }$ and the set $G_{ad,\chi
^{-1}}^{\wedge }$ contain equivalent representations. We write $(F^{\times
})^{\wedge }/\sim$ for the set of characters modulo the equivalence
relation: $\chi \sim \chi ^{\prime }$ if and only if $\chi ^{\prime }=\chi
^{-1}$.

The unramified case contains the complementary series and special
representation as well, so will be handled later.

\subsubsection{Borel structure on $G_{u,\chi }^{\wedge }\;,\chi \in
A_u^{\wedge }\;{\rm ramified}$}

The Borel structure on $G_{u,\chi }^{\wedge },$ when $\chi $ does not belong
to the orbit of some $sign_\tau $ under $X_{un}(A)$ is Borel isomorphic to
the Borel structure on the circle. The Borel structure on $G_{u,\chi
}^{\wedge },$ when $\chi $ {\em does} belong to the orbit of some $sign_\tau 
$ under $X_{un}(A)$ is Borel isomorphic to the Borel structure on the
half-circle union that on the set consisting of two (closed) points.

\subsubsection{Borel structure on $G_{ad,\chi }^{\wedge },\;\chi \in
A^{\wedge }\;{\rm ramified}$}

The Borel structure on $G_{ad,\chi }^{\wedge },$ when $\chi $ does not
belong to the orbit of some $sign_\tau $ under $X_{un}(A)$ is Borel
isomorphic to the Borel structure on $\{{\bf C}^{\times }\}/\{\pm 1\}.$
The Borel structure on $G_{ad,\chi }^{\wedge },$ when $\chi $ {\em does}
belong to the orbit of some $sign_\tau $ under $X_{un}(A)$ is Borel
isomorphic to the Borel structure on ${\bf C}^{\times }/\{\pm 1\}$ with two
double points.

\subsection{Complementary series}

This is sometimes called the supplementary series. For $V$ as above, $%
-1<\rho <1,\;\rho \neq 0,$ define $T_\rho :G\rightarrow End(V)$ by $T_\rho
=T_\chi ,\;$where $\chi (x)=|x|^{-\rho },$ for $x\in F^{\times }.$

{\bf Lemma} (Gelfand-Graev):{\em \ Let }$-1<\rho <1,\;\rho \neq 0.$

{\em (1) If }$\rho _1\neq \pm \rho _2${\em \ then }$T_{\rho _1}${\em \ is
not equivalent to }$T_{\rho _2},$

{\em (2) }$T_\rho \cong T_{-\rho },$

{\em (3) }$T_\rho ${\em \ is unitary,}

{\em (4) for }$0<\rho <1,\;T_\rho ${\em \ is irreducible}$.$

The irreducible representations in (4) are called {\bf complementary series
representations}. We shall see later that these representations are
``non-tempered''.

\subsection{The special representation and the trivial representation}

For $V$ and $T_\chi $ as above, we have

{\bf Lemma} (Gelfand-Graev):{\em \ Let }$\chi (x)=|x|^{-1},$ for $x\in
F^{\times }.$

(1) $T_\chi $ is reducible,

(2) there is an irreducible, infinite-dimensional, unitary,
square-integrable subrepresentation $T_{sp}$ of $T_\chi ,$

(3) The quotient of $T_\chi $ by $T_{sp}$ is 1-dimensional and unitary.

The representation in (2) is called the {\bf special representation}. The
representation in (3) is called the {\bf trivial representation}. We will
see later that the special representation is ``tempered'' but the trivial
representation is not.

\subsection{The $G_\sigma ^{\wedge }$ in the unramified case}

We define 
\[
G_{u,\chi }^{\wedge }=\{T_\rho ,\;T_{\chi ^{\prime }},\;T_{sign_\epsilon
}^{+},\;T_{sign_\epsilon }^{-}\;,1,T_{sp\;}|\; 
\begin{array}{c}
\begin{array}{c}
\chi ^{\prime }(x)=z^{\upsilon _F(x)},\; \\ 
z=e^{i\theta },\;\;0\leq \theta <\pi ,
\end{array}
\\ 
0<\rho <1,
\end{array}
\;\}, 
\]
if $\chi $ is an unramified unitary character of $A$,\ and, if $\chi $ is
unramified on $A$ 
\[
G_{ad,\chi }^{\wedge }=\{T_{\chi ^{\prime }},\;T_{sign_\epsilon
}^{+},\;T_{sign_\epsilon }^{-},1,T_{sp\;}\;|\; 
\begin{array}{c}
\chi ^{\prime }(x)=z^{\upsilon _F(x)}\neq |...|^{-1},\;sign_\varepsilon , \\ 
z\in {\bf C}^{\times }/\{\pm 1\},
\end{array}
\;\}, 
\]
where $\{\pm 1\}$ acts on ${\bf C}^{\times }$ by means of inversion. Note
that the set $G_{u,\chi }^{\wedge }$ and the set $G_{u,\chi ^{-1}}^{\wedge }$
(resp., the sets $G_{ad,\chi }^{\wedge }$ and $G_{ad,\chi ^{-1}}^{\wedge }$)
contain equivalent representations.

\subsubsection{Borel structure on $G_{u,\chi }^{\wedge },$ $\chi \in
A_u^{\wedge }$\ unramified}

This part of the dual consists of a circle, with opposing points identified
(except the point $z=1$ is not identified with $z=-1,$ since $z=-1$
corresponds to the double point) , union a set of four points, union the
line segment $(0,1)$. On each connected component, the Borel structure can
be taken to be that inherited from the half-circle, the segment $(0,1),$ or
the singleton set.

\subsubsection{Borel structure on $G_{ad,\chi }^{\wedge },\;\chi \in
A^{\wedge }$ unramified}

This part of the dual consists of the ${\bf C}^{\times }/\sim$ with
one double point, union a set of two (closed) points. On each connected
component, the Borel structure can be taken to be that inherited from
${\bf C}^{\times }/\{\pm 1\}$ or the singleton set.

\subsection{Cuspidal representations}

For each quadratic extension $E/F$ there is an infinite series of
irreducible cuspidal representations.

Let $\chi \in (E^{\times 1})^{\wedge }$, where $E^{\times 1}$ denotes the
kernel of the norm map, and define $T_\chi :G\rightarrow L^2(G)$ by 
\[
(T_\chi (g)f)(x)=\int_FK_\chi (g\;|\;x,y)f(y)dy, 
\]
where $K_\chi (g\;|\;x,y)$ is a certain distribution-valued kernel function
[GG]. Let $\chi _0$ denote the unique character of $E^{\times 1}$ of order $%
2 $.

{\bf Lemma} (Gelfand-Graev,Sally-Shalika) {\em Let }$\chi ,\chi _1,\chi
_2\in (E^{\times 1})^{\wedge }.$

{\em (1) For }$\chi \neq \chi _0,${\em \ }$T_\chi ${\em \ decomposes into a
direct sum of 2 irreducible, inequivalent subrepresentations }$T_\chi
^{+},\;T_\chi ^{-},$

{\em (2) For }$\chi =\chi _0,${\em \ }$T_\chi ${\em \ decomposes into a
direct sum of 4 irreducible, inequivalent subrepresentations }$T_\chi
^1,\;T_\chi ^2,\;T_\chi ^3,\;T_\chi ^4,$

{\em (3) For }$\chi \neq \chi _0,${\em \ we have }$T_\chi ^{+}\cong T_{\chi
^{-1}}^{+},\;T_\chi ^{-}\cong T_{\chi ^{-1}}^{-},$

{\em (4) If }$\chi \neq \chi _0${\em \ and }$\chi _1\neq \chi _2,\chi
_2^{-1} ${\em \ then }$T_{\chi _1}^{+}${\em \ is not equivalent to }$T_{\chi
_2}^{+}${\em \ and }$T_{\chi _1}^{-}${\em \ is not equivalent to }$T_{\chi
_2}^{-},$

{\em (5) If }$\chi _1,\chi _2\neq \chi _0${\em \ then }$T_{\chi _1}^{+}${\em %
\ is not equivalent to }$T_{\chi _2}^{-},$

{\em (6) For each cuspidal representation }$\sigma ${\em \ of }$G,${\em %
there is a quadratic extension }$E/F${\em \ and a character }$\chi \in
(E^{\times 1})^{\wedge }${\em \ such that }$\sigma =T_\chi ^\alpha ,${\em \
for some }$\alpha \in \{+,-,1,2,3,4\}.$

The treatment of the subrepresentations $T_\chi ^1,\;T_\chi ^2,\;T_\chi
^3,\;T_\chi ^4$ given in [GG] was not correct but was rectified in [SS1].

If $\sigma \in G_c^{\wedge }=G_{u,c}^{\wedge }$ then $G_{ad,\sigma }^{\wedge
}=G_{u,\sigma }^{\wedge }=\{\sigma \}$. The set of course gets the discrete
topology.

\subsection{Summary}

The admissible dual $G_{ad}^{\wedge }$ is a countable union of copies of$\;%
{\bf C}^{\times }/\{\pm 1\},$ union a ${\bf C}^{\times }/\{\pm 1\}$ with two
double points, union a ${\bf C}^{\times }/\{\pm 1\}$ with one double point
and two (closed) points, union a countable number of singleton sets.

The unitary dual $G_u^{\wedge }$ is a countable union of circles, union a
``half-circle'' with two double points, union a ``half-circle'' with one
double point and two (closed) points, union the line segment $(0,1)$ mod $%
2\pi /\ln (q)$, union a countable number of singleton sets. 


\section{The Fourier transform}

For $f\in C_c^\infty (G)$ and $\pi \in \Pi \in G_{ad}^{\wedge }$, we define
the {\bf operator-valued Fourier transform} of $f$ at $\pi $ by 
\[
\pi (f)=\int_Gf(x)\pi (x)\;dx,\;\;\;f\in C_c^\infty (G). 
\]

{\bf Lemma}: {\em If }$f,g\in C_c^\infty (G)${\em \ then }$\pi (f*g)=\pi
(f)\pi (g).$

Let $G^{\prime }=\{g\in G\;|\;g\;{\rm has\;distinct\;eigenvalues}\},$ called
the {\bf regular set} of $G.$ In particular, every element in the regular
set is semi-simple (i.e., diagonalizable).

Harish-Chandra showed that there exists a locally constant function, denoted 
$\Theta _\pi $ and called the {\bf character} of $\pi ,$ on the regular set $%
G^{\prime }$ which represents the trace of $\pi :$%
\[
trace(\pi (f))=\int_G\Theta _\pi (x)f(x)dx, 
\]
for all $f\in C_c^\infty (G).$ This trace, which we sometimes denote simply
by $\Theta _\pi (f)$, is called the {\bf Fourier transform} of $f$ at $\pi .$

\subsection{The Schwartz space and the tempered dual}

Let $||g||=\max (|a|,|b|,|c|,|d|),$ where $g=\left[ 
\begin{array}{cc}
a & b \\ 
c & d
\end{array}
\right] \in G,$ and let $\sigma (g)=\log ||g||.$ For each compact open
subgroup $K\subset \subset G$, let 
\[
{\cal C}_K(G)=\{f\in C_c(G//K)\;|\; 
\begin{array}{c}
|f(x)|<<_{r,f}\frac{\Xi (x)}{(1+\sigma (x))^r},\;\;\forall x\in G, \\ 
{\rm for\;each}\;r>0
\end{array}
\}, 
\]
where $C_c(G//K)$ denotes the space of compacty supported functions which
are bi-$K$-invariant, 
\[
\Xi (x)=\int_{K_0}\delta _B(xk)^{-1/2}dk. 
\]
Here, for $x\in G^{\prime }$, we have $\delta _B(x)=|\det (Ad(x_d))_{{\frak n%
}}|,$ where $x_d$ denotes a diagonalization of $x$ in $SL(2,\overline{F}),$
where $\overline{F}$ denotes a separable algebraic closure of $F$ and the
valuation $|...|$ has been extended to $\overline{F}$. It's known that there
are constants $c_1>0,\;c_2>0,\;N>0$ such that 
\[
c_1\leq \Xi (a)\leq c_2(1+\sigma (a))^N, 
\]
for all $a\in A^{+}=\{a\in A\;|\;a=\left[ 
\begin{array}{cc}
u\pi ^n & 0 \\ 
0 & u^{-1}\pi ^{-n}
\end{array}
\right] ,\;u\in O_F^{\times },\;n\geq 0\}.$ We topologize ${\cal C}_K(G)$
via the semi-norms 
\[
\upsilon _n(f)=\sup_{x\in G}|f(x)|\frac{(1+\sigma (x))^n}{\Xi (x)}. 
\]
Let 
\[
{\cal C}(G)=\bigcup_K{\cal C}_K(G), 
\]
where $K$ runs over all compact open subgroups of $G$. This is the {\bf %
Schwartz space} of $G$. Let $S$ denote the collection of all seminorms on $%
{\cal C}(G)$ whose restriction to each ${\cal C}_K(G)$ is continuous. In the
semi-norm topology induced by $S$, the Schwartz space is a complete locally
convex topological vector space.

{\bf Lemma} {\em (1) }${\cal C}(G)\subset L^2(G),$

{\em (2) }${\cal C}(G)${\em \ is an algebra under convolution.}

We call a distribution $D$ on $C_c^\infty (G)$ {\bf tempered} if it extends
continuously to ${\cal C}(G).$ We call an admissible representation $\pi $ 
{\bf tempered} if the distribution $f\longmapsto \Theta _\pi (f)$ extends
continuously to ${\cal C}(G).$ Let $G_t^{\wedge }$ denote the set of
equivalence classes of irreducible tempered representations of $G$, the {\bf %
tempered dual}.

{\bf Theorem} (Harish-Chandra): $G_t^{\wedge }\subset G_u^{\wedge }.$

In fact, for our case of $SL(2),$ we have 
\[
G_t^{\wedge }=G_u^{\wedge }-\{1,T_\rho \;|\;0<\rho <1\}. 
\]

Let 
\[
\begin{array}{c}
G_{t,\sigma }^{\wedge }=\{{\rm all\;irreducible\;inequivalent\;tempered\;}
\\ 
\;\;\;\;\;\;\;\;\;\;\;\;\;\;\;\;\;\;\;\;\;\;\;\;\;\;\;\;\;{\rm %
constituents\;of\;}I_M(\chi \sigma )\;\;|\;\;\chi \in X_{un}(M)\},
\end{array}
\]

We shall see that the following holds

{\bf Theorem:} $G_t^{\wedge }=\bigcup_{M\in {\cal L}(G)}\bigcup_{\sigma \in
M_{u,c}^{\wedge }}G_{t,\sigma }^{\wedge },$

We will see that $G_{t,\sigma }^{\wedge }=G_{u,\sigma }^{\wedge }$ if $%
\sigma\in\hat{A}_u$ is ramified. If $\sigma$ is unramified then $G_{t,\sigma
}^{\wedge }$ is the same as $G_{u}^{\wedge }$, except that it doesn't
contain the $T_\rho$ or the trivial representation. In other words, the only
non-tempered unitary representations are the trivial representation and the
complementary series. This will follow from the character identities
discussed next.

\subsection{Character calculations}

Let $S$ denote a complete set of representatives of conjugacy classes of
Cartan subgroups of $G$.\ We may and do assume that $A\in S.$ For each $T\in
S$, let $W_T=N_G(T)/T$ denote the Weyl group of $T$.

We begin with the

{\bf Lemma} (Weyl integration formula [S, p. 198]) {\em For each }$f\in
C_c^\infty (G),$%
\[
\int_Gf(g)=\sum_{T\in S}\frac 1{\#[W_T]}\int_T|D(t)|^{1/2}F_f^T(t)dt, 
\]

{\em where} 
\[
F_f^T(t)=|D(t)|^{1/2}\int_{T\backslash G}f(x^{-1}tx)\frac{dx}{dt}, 
\]

{\em and where }$D(t)=D_{G/T}(t)=\det (1-Ad(t))_{{\frak g}/{\frak t}%
}.\;Here\;{\frak g},\;{\frak t}\;denote\;the\;Lie\;algebras\;of\;G,\;T.$

Note that $F_f^T(t)=F_f^T(t^w)$, where $w\in W_T,$ so that $%
F_f^A(a)=F_f^A(a^{-1}),$ since $a^w=a^{-1},$ where $w=\left[ 
\begin{array}{cc}
0 & -1 \\ 
1 & 0
\end{array}
\right] .$

\subsubsection{Principal series}

Note that 
\[
\begin{array}{c}
D_{G/A}(\left[ 
\begin{array}{cc}
a & 0 \\ 
0 & a^{-1}
\end{array}
\right] )=\det (1-Ad(\left[ 
\begin{array}{cc}
a & 0 \\ 
0 & a^{-1}
\end{array}
\right] ))_{{\frak g}/{\frak a}} \\ 
=(1-a^2)(1-a^{-2})=-(a-a^{-1})^2.
\end{array}
\]

{\bf Lemma} {\em (1) The map }$f\longmapsto F_f^A${\em \ defines a
surjection }$C_c^\infty (G)\rightarrow C_c^\infty (A).$

{\em (2) The map }$f\longmapsto F_f^A${\em \ defines a continuous surjection 
}${\cal C}(G)\rightarrow {\cal C}(A).$

{\bf proof:} (1) For $f\in C_c^\infty (G)$ and $a\in A^{\prime }=A\cap
G^{\prime },$ the function $x\longmapsto f(x^{-1}ax)$ is compactly supported
on $A\backslash G$ and 
\[
\begin{array}{c}
F_f^A(a)=|D(a)|^{1/2}\int_{A\backslash G}f(x^{-1}ax) \frac{dx}{da} \\ 
=\delta _B(a)^{1/2}\int_N\int_{K_0}f(kank^{-1})dkdn \\ 
=\delta _B(a)^{1/2}\int_N\overline{f}(an)dn,
\end{array}
\]
where 
\[
\overline{f}(x)=\int_{K_0}f(kxk^{-1})dk. 
\]

Since the function $x\longmapsto f(x^{-1}ax)$ is compactly supported on $%
A\backslash G$ the function $x\longmapsto \overline{f}(x^{-1}ax)$ is
compactly supported on $A\backslash G/K_0.$ We claim that given any compact
set $C\subset \subset G,$ there is an $a\in A$ with sufficiently large $%
\sigma (a)$ such that $an\notin C,$ for all $n\in N.$ Indeed, we may map $C$
onto $\overline{C}\subset G/K_0$ via $G\rightarrow G/K_0$ and $aN$ onto $%
\overline{aN}$. Since $G/K_0$ is non-compact but $\overline{C}$ is compact
in $G/K_0$, the Iwasawa decomposition implies that we may find an $a\in A$
such that $\overline{aN}\cap \overline{C}=\emptyset $. From the claim it
follows that $F_f^A(a)$ is compactly supported on $A^{\prime }=A\cap
G^{\prime }.$

Moreover, since $f$ is locally constant so is $\overline{f}$. Since $%
\overline{f}$ is compactly supported, $\int_N\overline{f}(an)dn$ is a
locally constant function on $A^{\prime }$.

It remains to show that $F_f^A(a)$ extends to a locally constant function on
all of $A$. This follows from either the Shalika germ expansion. It is also
an immediate corollary of part (2), which we turn to next.

(2) The existence of the orbital integral $F_f^A(a)$ follows from Clozel. By
Harish-Chandra, the map $f\longmapsto f^B$ defines a continuous map ${\cal C}%
(G)\rightarrow {\cal C}(A),$ where 
\[
f^B(a)=\delta _B(a)^{1/2}\int_Nf(an)dn 
\]

Note that both $\Xi $ and $\sigma $ are bi-$K_0$-invariant, so the map $%
f\longmapsto \overline{f}$ defines a continuous map ${\cal C}(G)\rightarrow 
{\cal C}(G).$ Consequently, the composition $f\longmapsto \overline{f}%
\longmapsto F_{\overline{f}}^A$ defines a continuous map ${\cal C}%
(G)\rightarrow {\cal C}(A).$ This completes the proof. $\Box $

Let $\chi \in (F^{\times })^{\wedge }$ and suppose that $\chi \neq \chi
^{-1}.$

{\bf Theorem} (Gelfand-Graev) {\em The character of the irreducible
representation }

$\pi =T_\chi ${\em \ is given on the regular set by } 
\[
\Theta _\pi (x)=\left\{ 
\begin{array}{c}
\frac{\chi (a)+\chi (a)^{-1}}{|a-a^{-1}|},\;\;x=\left[ 
\begin{array}{cc}
a & 0 \\ 
0 & a^{-1}
\end{array}
\right] ^g,\; {\rm for\;some}\;g\in G,\;a\neq \pm 1, \\ 
0,\;\;\;\;\;\;\;\;\;\;\;\;\;\;\;\;\;\;\;\;\;\;\;\;\;\;\;\;\;{\rm \;if\;}%
\;\;\;x\in G^{\prime }-A^G,
\end{array}
\right. 
\]

{\em where }$A^G=\{a^g\;|\;a\in A,\;g\in G\}.${\em \ }

This and the Weyl integration formula imply that 
\[
\begin{array}{c}
\Theta _\pi (f)=\int_G\Theta _\pi (x)f(x)dx \\ 
=\frac 12\int_{F^{\times }}(\chi (a)+\chi (a)^{-1})F_f^A( 
\begin{array}{cc}
a & 0 \\ 
0 & a^{-1}
\end{array}
)d^{\times }a \\ 
=(F_f^A)^{\wedge }(\chi ),
\end{array}
\]
where 
\[
\begin{array}{c}
\varphi ^{\wedge }(\omega )=\int_{F^{\times }}\omega (x)\varphi (x)d^{\times
}x \\ 
=\sum_{n=-\infty }^\infty \omega (\pi )^n\int_{O_F{}^{\times }}\omega
(u)\varphi (\pi ^nu)du,\ \ \ \ \ \ \omega \in (F^{\times })_u^{\wedge }.
\end{array}
\]
denotes the Fourier transform on $F^{\times }.$ Note that this formula for
the character extends to the Schwartz space provided $\chi $ is unitary. In
fact, we have

{\bf Lemma}: {\em A principal series representation }$T_\chi ${\em \ is
tempered if and only }$\chi ${\em \ is unitary.}

The above formulas allows us to classify the image of the Fourier transforms
of a {\em unitary} principal series representation on either $C_c^\infty (G)$
or the Schwartz space.

First, suppose $f\in C_c^\infty (G)$, so $F_f^A\in C_c^\infty (A).$

{\bf Lemma}:{\em \ For }$\varphi \in C_c^\infty (F^{\times }),${\em \ the
image }$C_c^\infty (F^{\times })^{\wedge }${\em \ of the Fourier transform }$%
\varphi \longmapsto \varphi ^{\wedge }(\omega )${\em \ is given by} 
\[
C_c^\infty (F^{\times })^{\wedge }=\{h\in C_c((F^{\times })_u^{\wedge
})\;|\; 
\begin{array}{c}
h\;is\;a\;trig\;polynomial\;on \\ 
each\;circle\;in\;(F^{\times })_u^{\wedge }
\end{array}
\;\}. 
\]

We omit the proof. (This lemma, and its Schwartz space analog below, is
proven using the inverse Fourier transform and the inversion formula, which
may be found in [Ta, pp. 43-44], [Ba, ch 12].) As a consequence of the above
lemma, since the map $C_c^\infty (G)\rightarrow C_c^\infty (A)$ is
surjective, it follows that the same description holds for the space of all
functions of $\chi $ of the form $(F_f^A)^{\wedge }(\chi ):$

{\bf Lemma}:{\em \ For }$f\in C_c^\infty (G),${\em \ the image }$C_c^\infty
(G)_{ps}^{\wedge }${\em \ of the Fourier transform }$f\longmapsto \Theta
_\pi (f),\;\pi =T_\chi ,\;\chi \in A_u^{\wedge }${\em \ is given by} 
\[
C_c^\infty (G)_{ps}^{\wedge }=\{h\in C_c(A_{u,c}^{\wedge })\;|\; 
\begin{array}{c}
h\;is\;a\;trig\;polynomial\;on \\ 
each\;circle\;in\;A_{u,c}^{\wedge }
\end{array}
\;\}. 
\]

Next, suppose $f\in {\cal C}(G)$, so $F_f^A\in {\cal C}(A).$

{\bf Lemma}: {\em For }$\varphi \in {\cal C}(F^{\times }),${\em \ the image }%
${\cal C}(F^{\times })^{\wedge }${\em \ of the Fourier transform }$\varphi
\longmapsto \varphi ^{\wedge }(\omega )${\em \ is given by} 
\[
{\cal C}(F^{\times })^{\wedge }=\{h\in C_c((F^{\times })_u^{\wedge })\;|\; 
\begin{array}{c}
\;h\;is\;a\;trig\;series\;on\;each\;circle \\ 
in\;(F^{\times })_u^{\wedge }\;which\;converges\;absolutely\; \\ 
along\;with\;all\;of\;its\;derivatives
\end{array}
\}. 
\]
Consequently, since the map ${\cal C}(G)\rightarrow {\cal C}(A)$ is
surjective, it follows that the same description holds for the space of all
functions of $\chi $ of the form $(F_f^A)^{\wedge }(\chi ):$

{\bf Lemma}: {\em For }$f\in {\cal C}(G),${\em \ the image }${\cal C}%
(G)_{ps}^{\wedge }${\em \ of the Fourier transform }$f\longmapsto \Theta
_\pi (f),\;\pi =T_\chi ${\em \ is given by} 
\[
{\cal C}(G)_{ps}^{\wedge }=\{h\in C_c(A_{u,c}^{\wedge })\;|\; 
\begin{array}{c}
\;h\;is\;a\;trig\;series\;on\;each\;circle \\ 
in\;A_{u,c}^{\wedge }\;which\;converges\;absolutely\; \\ 
along\;with\;all\;of\;its\;derivatives
\end{array}
\}. 
\]

\subsubsection{Reducible principal series}

Let $\chi \in (F^{\times })^{\wedge }$ and suppose that $\chi =sign_\tau ,$
for some $\tau =\varepsilon ,\pi ,\varepsilon \pi .$ Let $\pi =T_\chi ^{+}$%
{\em \ or }$\pi =T_\chi ^{-}$. Following [SS] or [Fr] (see also [A]), the
character $\Theta _\pi $ of a reducible principal series representation is
given on the regular set by 
\[
\Theta _\pi (x)=\left\{ 
\begin{array}{c}
\frac{sign_\tau (x)}{|a-a^{-1}|},\;\;\;\;x=\left[ 
\begin{array}{cc}
a & 0 \\ 
0 & a^{-1}
\end{array}
\right] ^g,\;for\;some\;g\in G,\;a\neq \pm 1,\; \\ 
e_{\pi ,\tau }(t)|D_T(t)|^{-1/2},\;\;\;\;\;\;x\;conjugate\;to\;t\in (T_\tau
\cup T_\tau ^{\#})\cap G^{\prime }, \\ 
0,\;\;\;\;\;\;\;\;\;\;\;\;\;\;\;\;\;\;\;\;\;\;\;\;\;\;\;\;\;\;\;\;\;\;\;\;\;%
\;\;\;\;\;\;\;\;\;otherwise,
\end{array}
\right. 
\]
where $T_\tau ,T_\tau ^{\#}$ are two non-split Cartans (which are not
conjugate if $-1\in (F^{\times })^2$ or if $\tau =\varepsilon )$ and where $%
e_{\pi ,\tau }(t)$ is a bounded locally constant function on $G^{\prime
}-A^G $. Using this, the Weyl integration formula, and the reasoning above,
it can be shown that

{\bf Lemma} {\em If }$\pi =T_\chi ^{+}${\em \ or }$\pi =T_\chi ^{-}${\em \
then }$\pi ${\em \ is tempered.}

\subsubsection{Complementary series}

Let $\chi =|...|^{-\rho }$ for some $0<\rho <1.$ As in the case of the
principal series, we see that $T_\rho $ is not tempered because $\chi $ is
not a unitary character.

{\bf Lemma} {\em If }$\pi =T_\rho ${\em \ then }$\pi ${\em \ is non-tempered.%
}

The character $\Theta _\pi $ of a reducible principal series representation
is given by the same formula as in the case of the principal series. This
and the Weyl integration formula imply that 
\[
\begin{array}{c}
\Theta _\pi (f)=\int_G\Theta _\pi (x)f(x)dx \\ 
=\frac 12\int_{F^{\times }}(|a|^{-\rho }+|a|^\rho )F_f^A( 
\begin{array}{cc}
a & 0 \\ 
0 & a^{-1}
\end{array}
)d^{\times }a \\ 
=M(F_f^A)(\rho ),
\end{array}
\]
where $M$ denotes the {\bf Mellin transform} 
\[
\begin{array}{c}
M\varphi (\rho )=\int_{F^{\times }}|x|^{-\rho }\varphi (x)d^{\times }x \\ 
=\sum_{n=-\infty }^\infty q^{\rho n}\int_{O_F{}^{\times }}\varphi (\pi
^nu)du.
\end{array}
\]

{\bf Lemma}:{\em \ For }$\varphi \in C_c^\infty (F^{\times }),${\em \ the
image }$MC_c^\infty (F^{\times })${\em \ of the\ Fourier transform }$\varphi
\longmapsto \Theta _\pi (f),\;\pi =T_\rho ${\em \ is given by} 
\[
MC_c^\infty (G)=\{h\in C^\infty
(0,1)\;|\;h\;is\;a\;finite\;Laurant\;series\;in\;q^\rho \;\}. 
\]

Consequently, since the map $C_c^\infty (G)\rightarrow C_c^\infty (A)$ is
surjective, it follows that the same description holds for the space of all
functions of $\chi $ of the form $M(F_f^A)(\rho ):$

{\bf Lemma}:{\em \ For }$f\in C_c^\infty (G),${\em \ the image }$MC_c^\infty
(G)${\em \ of the Mellin transform }$f\longmapsto M(F_f^A)(\rho )${\em \ is
given by} 
\[
C_c^\infty (F^{\times })^{\wedge }=\{h\in C^\infty
(0,1)\;|\;h\;is\;a\;finite\;Laurant\;series\;in\;q^\rho \;\}. 
\]

\subsubsection{The special representation}

Let $\pi =T_{sp}$. Then{\em \ } 
\[
\Theta _\pi (x)=\left\{ 
\begin{array}{c}
\frac{|a|+|a|^{-1}}{|a-a^{-1}|}-1,\;\;x=\left[ 
\begin{array}{cc}
a & 0 \\ 
0 & a^{-1}
\end{array}
\right] ^g,\;for\;some\;g\in G,\;a\neq \pm 1, \\ 
-1,\;\;\;\;\;\;\;\;\;\;\;\;\;\;x\in G^{\prime }-A^G.
\end{array}
\right. 
\]
It follows from this and the Weyl integration formula that 
\[
\begin{array}{c}
\Theta _\pi (f)=\int_G\Theta _\pi (x)f(x)dx \\ 
=\frac 12\int_{F^{\times }}(|a|+|a|^{-1}-|a-a^{-1}|)F_f^A( 
\begin{array}{cc}
a & 0 \\ 
0 & a^{-1}
\end{array}
)d^{\times }a \\ 
-\sum_{T\in S-\{A\}}\frac 12\int_T|D(t)|^{1/2}F_f^T(t)dt.
\end{array}
\]
This converges even for $f$ in the Schwartz space, so we have

{\bf Lemma}: {\em If }$\pi =T_{sp}${\em \ then }$\pi ${\em \ is tempered.}

\subsubsection{The trivial representation}

Let $\pi =1$ so{\em \ } 
\[
\Theta _\pi (x)=1. 
\]
To see that 
\[
\Theta _\pi (f)=\int_Gf(x)dx 
\]
diverges in general on the Schwartz space, we use [S, Lemma 4.2.5] to
compute 
\[
\begin{array}{c}
\int_G \frac{\Xi (x)}{(1+\sigma (x))^r}dx=\sum_{a\in A^{+}/A^0}\mu (K_0aK_0)%
\frac{\Xi (a)}{(1+\sigma (a))^r} \\ 
\approx \sum_{a\in A^{+}/A^0}\mu (K_0aK_0)\delta _B(a)^{1/2} \frac{\delta
_B(a)^{-1/2}\Xi (a)}{(1+\sigma (a))^r} \\ 
\approx \sum_{a\in A^{+}/A^0}\delta _B(a)^{1/2}\frac 1{(1+\sigma
(a))^{r-r_0}},
\end{array}
\]
for some constant $r_0$ and all $r>0$. (Here $\mu $ denotes the volume.)
This diverges for all $r$, so the trivial representation is not tempered.
(Here $f(a)\approx g(a)$ means $f(a)<<g(a)<<f(a)$.)

{\bf Lemma} {\em If }$\pi =1${\em \ then }$\pi ${\em \ is non-tempered.}

\subsubsection{The cuspidal representations}

The characters of the cuspidal representations will not be given here - we
refer to [SS2]. They are all of the form
\[
\Theta _\pi (x)=\left\{ 
\begin{array}{c}
\alpha (x)+ \frac{\beta (x)}{|D(x)|^{1/2}},\;\;x\in G^{\prime }-A^G, \\ 
0,\;\;\;\;\;\;\;\;\;\;\;\;\;\;x\notin G^{\prime }-A^G,
\end{array}
\right. 
\]
where $\alpha ,\beta $ are uniformly bounded locally constant functions on $%
G^{\prime }-A^G.$ From this and our previous estimates, it is not hard to
see that the cuspidal representations are tempered (see also [S, Theorem
4.5.1, Theorem 4.5.10]).

{\bf Lemma}: {\em If }$\pi \in G_c^{\wedge }${\em \ then }$\pi ${\em \ is
tempered.}

\section{The image of the Fourier transform}

\subsection{Admissible Paley-Wiener Theorem for smooth functions}

The following result classifies the image of $C_c^\infty (G)$ under the
Fourier transform:

{\bf Theorem} (Bernstein-Deligne-Kazhdan) {\em The image of the Fourier
transform }$\pi \longmapsto \Theta _\pi (f),${\em \ for }$f\in C_c^\infty
(G) ${\em \ and }$\pi \in \Pi \in G_{ad}^{\wedge }${\em , consists of the
space of functions }$h${\em \ on }$G_{ad}^{\wedge }${\em \ such that}

{\em (1) }$h${\em \ is supported on finitely many connected components of }$%
G_{ad}^{\wedge }${\em ,}

{\em (2) }$h${\em \ is regular on }${\bf C}^{\times }/\{\pm 1\}${\em \ in $%
G_{ad}^{\wedge }$\ , regarded as a complex algebraic variety.}

Let $C_c^\infty (G)_{ad}^{\wedge }$ denote the space of functions described
in the above theorem. We call it the {\bf admissible Paley-Wiener space}.

\subsection{Unitary Paley-Wiener Theorem for smooth functions}

We use the information above to prove the following

{\bf Theorem} {\em The image of the Fourier transform }$\pi \longmapsto
\Theta _\pi (f),${\em \ for }$f\in C_c^\infty (G)${\em \ and }$\pi \in \Pi
\in G_u^{\wedge }${\em , consists of the space of functions }$h${\em \ on }$%
G_u^{\wedge }${\em \ such that}

{\em (1) }$h${\em \ is supported on finitely many connected components of }$%
G_u^{\wedge }${\em ,}

{\em (2) }$h${\em \ is a trigonometric polynomial on each circle (or
half-circle) in }$G_u^{\wedge }${\em \ which it is supported on,}

{\em (3) if }$h${\em \ is supported on the complementary series then }$h$%
{\em \ is a finite Laurant series in }$x=q^\rho .$

{\bf proof}: (1) follows from the admissible Paley-Wiener theorem.

(2) follows from the section on characters of principal series
representations.

(3) follows from the section on characters of complementary series
representations.

This completes the proof since $h$ can be arbitrary on the discrete part of
the unitary dual. $\Box $

Let $C_c^\infty (G)_u^{\wedge }$ denote the space of functions described in
the above theorem. We call it the {\bf unitary Paley-Wiener space}.

\subsection{Tempered Paley-Wiener Theorem for Schwartz functions}

We also have the following

{\bf Theorem} {\em The image of the Fourier transform }$\pi \longmapsto
\Theta _\pi (f),${\em \ for }$f\in {\cal C}(G)${\em \ and }$\pi \in \Pi \in
G_t^{\wedge }${\em , consists of the space of functions }$h${\em \ on }$%
G_t^{\wedge }${\em \ such that}

{\em (1) }$h${\em \ is supported on finitely many connected components of }$%
G_t^{\wedge }${\em ,}

{\em (2) }$h${\em \ is a trigonometric series on each circle (or
half-circle, ignoring the double point if it occurs) in }$G_t^{\wedge }${\em %
\ which it is supported on, converging absolutely along with all of its
derivatives.}

{\bf proof}: (1) follows from the admissible Paley-Wiener theorem.

(2) follows from the section on characters of principal series
representations.

This completes the proof since $h$ can be arbitrary on the discrete part of
the tempered dual. $\Box $

Let ${\cal C}(G)^{\wedge }$ denote the space of functions described in the
above theorem. We call it the {\bf tempered Paley-Wiener space}.

\section{Linear functionals on the tempered Paley-Wiener space}

\subsection{A result of L. Schwartz}

The space of linear functionals on the tempered Paley-Wiener space can be
successfully described using the theory of distributions on the real line:

{\bf Theorem} (L. Schwartz [Sch, ch. III, Th\'eor\`eme XXI]) {\em Let }$S^1$%
{\em \ denote the unit circle. If }$T\in C^\infty (S^1)^{\prime }${\em \
then there is a continuous function }$u${\em \ on }$S^1${\em \ and an
integer }$n\geq 0${\em \ such that }$T=\frac{d^n}{d\theta ^n}u${\em \ (as
distributions).}

As a corollary of this and the above Paley-Wiener theorem, we have the
following

{\bf Proposition}: {\em The dual space of the tempered Paley-Wiener space } $%
{\cal C}(G)^{\wedge }$ {\em \ is isomorphic to the direct product} 
\[
\bigoplus_{d\in \pi _0(A_{u,c}^{\wedge })}\bigcup_{n\geq
0}C(S_d)^{(n)}\bigoplus \bigoplus_{\sigma \in \pi _0(G_{lds}^{\wedge })}{\bf %
C}\delta _\sigma , 
\]
{\em where }$\pi _0(X)${\em \ denotes the set of connected components of a
topological space }$X,$ {\em where} $G_{lds}^{\wedge }$ {\em denotes the
union of the reducible principal series, the special representation, and the
cuspidal representations,}$\;{\em and\;where\;}\delta _\sigma $ {\em denotes
the evaluation functional on the space of (constant) functions on the
(singleton) elements of } $\pi _0(G_{lds}^{\wedge })${\em , and, if we write
the unitary dual }$A_u^{\wedge }=A_{u,c}^{\wedge }${\em \ as a union of
circles }$S_d\cong {\em {\bf R}/\frac{2\pi }{\log (q)}{\bf Z,\;}where}${\em %
\ } 
\[
C(S_d)^{(n)}=\{\frac{d^n}{d\theta ^n}u\;|\;u\;{\rm continuous\;on\;}S_d\}, 
\]

\subsection{The Fourier transform of a tempered distribution}

As a consequence of this, we can ``completely describe'' the Fourier
transform of a tempered distribution:

{\bf Theorem}: {\em If }$T${\em \ is a tempered distribution on }$G${\em \
then there is a } 
\[
T^{\wedge }\in \bigoplus_{d\in \pi _0(A_{u,c}^{\wedge })}\bigcup_{n\geq
0}C(S_d)^{(n)}\bigoplus \bigoplus_{\sigma \in \pi _0(G_{lds}^{\wedge })}{\bf %
C}\delta _\sigma 
\]
{\em \ such that} 
\[
T(f)=T^{\wedge }(f^{\wedge }), 
\]
{\em for all }$f\in {\cal C}(G)${\em , where }$f^{\wedge }(\pi )=\Theta _\pi
(f)${\em \ denotes the Fourier transform of }$f${\em \ at }$\pi \in \Pi \in
G_t^{\wedge }.$

\section{Linear functionals on the admissible Paley-Wiener space}

\subsection{Analytic functionals}

A large part of the space of linear functionals on the admissible
Paley-Wiener space can be described using the theory of analytic functionals
(see H\"ormander [H1, ch. IV], [H2, ch. IX]). Our description of the Fourier
transform of an admissible distribution is unfortunately much less explicit
than in the tempered case treated in the previous section. However, if the
admissible distribution has ``exponential growth'' then, as we will see,
more precise results can be given.

In spite of what we have seen already in the previous section, the following
may still be of interest

{\sc Question} Is there is an analogous description for the linear
functionals on the unitary Paley-Wiener space using the theory of (Sato)
hyperfunctions?

(This question is inspired by the relationship between analytic functionals
and hyperfunctions discussed in [H2, \S 9.2].)

We begin with some definitions:

Let $\Omega \subset G_\sigma ^{\wedge }\hookrightarrow X_{un}(M)\cong ({\bf C%
}^{\times })^n,\;\sigma \in M_c^{\wedge },\;$ some $n\geq 1,$ denote a
connected component of $G_{ad}^{\wedge }$, which we may regard as an complex
algebraic variety or as a complex analytic manifold. Let $R(X_{un}(M))$
denote the space of regular functions on $X_{un}(M)$, regarded as a variety,
and $A(X_{un}(M))$ denote the space of analytic functions on $X_{un}(M)$,
regarded as a manifold. We give $A(X_{un}(M))$ the topology of uniform
convergence on compacta, so that $R(X_{un}(M))\subset A(X_{un}(M))$ is
dense. Following H\"ormander ([H1, ch. IV]), an {\bf analytic functional} on 
$X_{un}(M)$ is an element of the topological dual of $A(X_{un}(M))$, denoted 
$A(X_{un}(M))^{\prime }$, where $A(X_{un}(M))$ is given the topology of
uniform convergence. We say that an analytic functional $\lambda $ is {\bf %
carried on} a compact set $K\subset \subset X_{un}(M)$ if for every
neighborhood $U\supset K$, there is a constant $C_U>0$ such that 
\[
|\lambda (f)|\leq C_U\sup _{z\in U}|f(z)|. 
\]
It follows from the definitions that for any $\lambda \in A(\Omega )^{\prime
}$, there exists a compact $K$ such that $\lambda $ is carried on $K$. The
analytic functionals carried on $K$ are in one-to-one correspondence, via
the Laplace transform, with those entire functions with exponential growth
(a result due to P\'olya- Ehrenpreis- Martineau, see [H1, Theorem 4.5.3]).

Let $PW(G)=C_c^\infty (G)_{ad}^{\wedge }$ denote the Paley-Wiener space and
let $PW(G)^{\prime }$ denote its algebraic dual space (as a vector space).
From the bijection $f\longmapsto f^{\wedge },\;C_c^\infty (G)\rightarrow
PW(G),$ we define the {\bf Fourier transform} of an invariant distribution $%
D $ on $C_c^\infty (G)$ to be the linear functional $D^{\wedge }\in
PW(G)^{\prime }$ such that 
\[
D(f)=D^{\wedge }(f^{\wedge }), 
\]
for all{\em \ }$f\in C_c^\infty (G).$ We say that the distribution $D$ is 
{\bf analytic} if the linear functional $D^{\wedge }$ restricted to $%
R(X_{un}(M))\subset PW(G)$ extends continuously an analytic functional on $%
A(\Omega )$.

{\bf Lemma}: {\em If }$D${\em \ is analytic then there is a compact }
$K\subset \subset \Omega ${\em \ such that }$D^{\wedge }${\em \ is carried on 
}$K${\em .}

Define a function $E_{M,a}$ on $X_{un}(M)$, by $E_{M,a}(I_M(\chi \sigma
))=\chi (a)\; $for\ each $\chi \in X_{un}(M),$ $a\in A_M$, where $A_M$
denotes the split component of the center of $M$ (in $M$). Let 
\[
H_K(a)=max_{\chi\in K}\log (||a||^{Re(s_\chi )}),\;\;\;\;\forall a\in A_M 
\]
where $s_\chi \in {\bf C}^n$ is an element representing $\chi \in X_{un}(M).$

{\bf Lemma} {\em The vector space span of all linear combinations of the
functions }$E_{M,a}${\em , }$a\in A_M${\em , is dense in }$A(X_{un}(M))$.

{\bf proof} The only cases we need to prove this are $M=G$ and $M=A$. The
statement is trivial if $M=G$ so assume $M=A$.

The trace Paley-Wiener theorem of Bernstein-Deligne-Kazhdan mentioned above
implies that each regular function $h$ on $X_{un}(A)$ may be realized as the
trace: there is an $f\in C_c^\infty (G)$ such that $h(\chi )=\Theta
_{I_A(\chi )}(f)$ for all $\chi \in X_{un}(A)$ such that the induced
representation $I_A(\chi )$ is irreducible. In this case, the character
formulas in the principal series case imply that 
\[
\begin{array}{c}
h(\chi )=\int_G\Theta _{I_A(\chi )}(g)f(g)dg=(F_f^A)^{\wedge }(\chi ) \\ 
=\int_AF_f^A(a)\chi (a)da=\sum_{a\in A,\ {\rm finite}} c_a\chi(a),
\end{array}
\]
for all $\chi \in X_{un}(A)$, where the sum over $a\in A$ runs over a finite
set (since $F_f^A(a)$ is compactly supported). From this the result easily
follows. $\Box $

Therefore to verify that an analytic functional carried on $K\subset \subset
X_{un}(M)$ it suffices to check that the condition $|\lambda (f)|\leq
C_U\sup _{z\in U}|f(z)|$ holds for all $f$ belonging to this span. Since 
\[
|E_{M,a}(\chi)|\leq ||a||^{Re(s_\chi)}, 
\]
this implies the following

{\bf Proposition}: {\em If }$D${\em \ is analytic and }$D^{\wedge }${\em \
is carried on a compact }$K\subset \subset X_{un}(M)${\em \ then for each }$%
\varepsilon >0${\em \ there is a constant }$C_\varepsilon >0${\em \ such that%
} 
\[
|D^{\wedge }(E_{M,a})|\leq C_\varepsilon \exp ((1+\varepsilon
)H_K(a)),\;\;\;\;\;\;\;\forall a\in A_M.\;\; 
\]

\subsection{The Fourier transform of an analytic distribution}

The result below describes the Fourier transform of an analytic distribution
on $G$.

{\bf Theorem}: {\em If }$D${\em \ is an analytic distribution then there is
  a linear functional }$D^{\wedge }\in PW(G)^{\prime }${\em \ such that}
$D(f)=D^{\wedge }(f^{\wedge }),${\em for all }$f\in C_c^\infty (G)${\em ,
  and, for each Levi }$M${\em \ of }$G${\em , and each }$\epsilon >0$
{\em there are constants }$C_\epsilon >0,\ C_M>0$, {\em such that} 
\[
|D^{\wedge }(E_{M,a})|\leq C_\varepsilon \exp ((1+\varepsilon )\sigma
(a)^{C_M}),\;\;\;\;\;\;\;\forall a\in A_M.\;\; 
\]
{\em Conversely, let }$D${\em \ be an invariant distribution such that its
Fourier transform }$D^{\wedge }\in PW(G)^{\prime }${\em \ satisfies the
  following property:} {\em \ for each Levi }$M${\em \ of }$G${\em , and each }
$\epsilon >0$ {\em there are constants }$C_\epsilon >0,\ C_M>0$, {\em such
that} 
\[
|D^{\wedge }(E_{M,a})|\leq C_\varepsilon \exp ((1+\varepsilon )\sigma
(a)^{C_M}),\;\;\;\;\;\;\;\forall a\in A_M.\;\; 
\]

{\em Then }$D${\em \ is analytic.}

\subsection{``Nice'' admissible distributions are analytic}

We shall define admissible distributions, which appear in the result below,
in the next section. Roughly speaking, the theorem below says that
admissible distributions of ``exponential growth'' on $SL(2)$, which are
very common but not well classified, are analytic.

{\bf Theorem}: {\em If }$D${\em \ is admissible and if there is a constant }$%
C>0$ {\em such that }$|D(a)|^{1/2}F_D(a)<<||a||^C${\em , for all }$a\in A$%
{\em \ then }${\em D}${\em \ is analytic.}

{\bf proof}: We need only check the condition of analyticity for the Levi
component $M=A$ since the case $M=G$ is essentially trivial due to the fact
that $A_M=\{\pm 1\}$ in that case. Furthermore, we need only prove the
estimate 
\[
|D^{\wedge }(E_{A,a})|\leq C_\varepsilon \exp ((1+\varepsilon )\sigma
(a)^{C_A}),\;\;\;\;\;\;\;\forall a\in A,\;\; 
\]
for ``sufficiently large'' $a\in A$ (that is, $a\in A$ for which $\sigma (a)$
is sufficiently large).

Fix an $a\in A$ and pick $f=f_a\in C_c^\infty (G)$ such that

(1) $f$ is supported in a neighborhood of the conjugacy class of $a$ in $G,$
and

(2)$\;F_f^A\in C_c^\infty (A)$ is supported in a neighborhood $U$ of $a$ in $%
A$ such that 
\[
\begin{array}{c}
(F_f^A)^{\wedge }(\chi )=\chi (a)+\chi (a)^{-1} \\ 
|F_f^A(x)|<<|a|+|a|^{-1},\;\;\;\;\forall x\in U.
\end{array}
\]
Since $D$ is admissible there is (by a result of Harish-Chandra stated
above) a function $F_D$ such that 
\[
D(f)=\int_GF_D(g)f(g)dg. 
\]
Therefore, by definition of $f=f_a$ , the definition of $E_{A,a}$, and the
Weyl integration formula, we have the estimate 
\[
\begin{array}{c}
D^{\wedge }(E_{A,a})|=2|D^{\wedge }(f^{\wedge })|=2|D(f)| \\ 
\leq 2\int_G|F_D(g)f(g)|dg \\ 
<<\sum_{T\in S,\ T\neq A } \frac
1{\#[W_T]}\int_T|D(t)|^{1/2}|F_D(t)F_f^T(t)|dt+%
\int_A|D(a)|^{1/2}|F_D(a)F_f^A(a)|da \\ 
<<1+\int_A|D(a)|^{1/2}|F_D(a)F_f^A(a)|da \\ 
<<1+|D(a)|^{1/2}(|a|+|a|^{-1})\max _{x\in U}|F_D(x)|<<||a||^{3+C}.
\end{array}
\]
(The fact that 
\[
\sum_{ T\in S,  \ T\neq A  }\frac
1{\#[W_T]}\int_T|D(t)|^{1/2}|F_D(t)F_f^T(t)|dt<<1 
\]
follows from the fact that the $T\neq A$ are compact, bounded in number, and 
$f$ is supported in a neighborhood of the conjugacy class of $a$ in $G$. If
we take $a$ ``sufficiently large'' in the sense above then this estimate
follows easily.)

This implies that $D$ is analytic by the previous theorem. $\Box $

{\bf Corollary}: {\em The space of all analytic distributions on }$G$ {\em %
includes the characters }$\Theta _\pi $, {\em for }$\pi \in \hat G_{ad}$.

{\bf proof}: There is a result of Harish-Chandra and Howe which implies that
the character of every irreducible admissible representation is an
admissible distribution. The bound $|D(a)|^{1/2}|F_D(a)|<<||a||^C$, for some 
$C=C(D)>0$, follows from the explicit character formulas above. $\Box $

\section{Admissible distributions}

\subsection{Definitions}

Let $K\subset \subset G$ be a compact open subgroup and for each irreducible
representation $\sigma $ of $K$, let $\chi _\sigma $ denote its character.

An invariant distribution $D$ on $G$ (not to be confused with the Weyl
discriminant function $D(x)$ on $G$) is called $(G,K)$-{\bf admissible at $%
\gamma \in G$} if for any compact open subgroup $K^{\prime }\subset K$ and
any irreducible representation $\sigma $ of $K^{\prime }$, we have $D*\chi
_\sigma =0$ on $\gamma K$ unless $G$ intertwines the trivial representation
on $K$ with $\sigma $. In other words, for each $f\in C_c^\infty (G)$ with $%
supp(f)\subset \gamma K$, we have $D(f*\chi _\sigma ^0)=0$ on $\gamma K$,
unless $G$ intertwines the trivial representation on $K$ with $\sigma $.
Here $\chi _\sigma ^0(x)$ denotes the extended character of $\sigma :$ $\chi
_\sigma ^0(x)=\chi _\sigma (x)$ if $x\in K^{\prime }$ and $\chi _\sigma
^0(x)=0$ if $x\in G-K^{\prime }$. If $U$ is a $G$-invariant open set in $G$
then we say that $D$ is {\bf admissible on }$U$ if it is $(G,K)$-admissible
at each $\gamma \in U$ in $G$ for some sufficiently small $K$ (which may
depend on $\gamma $). Note that it suffices to verify this for each regular $%
\gamma \in U\cap G^{\prime }.$

According to Harish-Chandra, if $D$ is admissible then there is a function $%
F_D$ on $G$ which is locally constant on the regular set $G^{\prime }$ such
that

(1) $|D(x)|^{1/2}F_D$ is locally bounded,

(2) for each $f\in C_c^\infty (G)$, 
\[
D(f)=\int_GF_D(x)f(x)dx. 
\]
In general, if $D$ is a distribution on $G$ which is given by an integral as
above, we say that $D$ is {\bf representable}. For representable
distributions, the $(G,K)$-admissibility condition at $\gamma $ becomes: for
each $f\in C_c^\infty (G)$ with $supp(f)\subset \gamma K$, we have 
\[
\begin{array}{c}
\int_GF_D(x)(f*\chi _\sigma ^0)(x)dx=\int_UF_D(x)\int_{K^{\prime
}}f(xk^{-1})\chi _\sigma ^0(k)dkdx \\ 
=\int_G(F_D*\chi _\sigma ^0)(x)f(x)dx=0,
\end{array}
\]
unless $G$ intertwines the trivial representation on $K$ with $\sigma $.
Therefore, the $(G,K)$-admissibility condition for represented distributions
is: $(F_D*\chi _\sigma ^0)(x)=0$ on $\gamma K,$ unless $G$ intertwines the
trivial representation on $K$ with $\sigma $.

\subsection{A descent to the Levis}

By the Weyl integration formula, we may write 
\[
D(f)=\sum_{T\in S}{\frac 1{|W_T|}}\int_T|D(t)|^{1/2}F_D(t)F_f^{G/T}(t)dt=%
\sum_MD_M(F_{\overline{f}^P}^{M/T}). 
\]
Moreover, here 
\[
D_M(h)=\sum_{T,\ M-elliptic } {\frac 1{|W_T|}}%
\int_T|D(t)|^{1/2}F_D(x)h(t)dt, 
\]
for all $h\in C_c^\infty (M),$ where the sum runs over all Cartans $T$ in $G$
which are $G$-conjugate to a Cartan of $M$ but not to any smaller $M^{\prime
}\in {\cal L}(G)$.

In the case of $SL(2)$ this amounts to saying that 
\[
D_G(f)=\int_GF_D(x)f(x)ch_{G-A^G}(x)dx, 
\]
and 
\[
D_A(f)=\int_GF_D(x)f(x)ch_{A^G}(x)dx=\frac
12\int_A|D(a)|^{1/2}F_D(a)F_f^A(a)da, 
\]
where $ch_S$ denotes the characteristic function of the set $S\subset G$.

{\bf Lemma}{\sc :} {\em If }$D${\em \ is }$G${\em -admissible and if }$M\in 
{\cal L}(G),M\neq G,${\em \ then }$D_M${\em \ is }$M${\em -admissible.}

{\sc Remark:} $D_G$ is admissible at $\gamma $\ provided 
\[
\begin{array}{c}
(ch_{G-A^G}\cdot F_D*\chi _\sigma ^0)(x)=0,\; {\rm on\;}\gamma K\;{\rm %
unless\;}G\;{\rm intertwines} \\ 
{\rm \;the\;trivial\;representation\;on\;}K\;{\rm with}\;\sigma .
\end{array}
\]
In this case, since $G-A^G$ is open in $G$, we see that for each $\gamma \in
G-A^G$ there is a compact open $K\subset \subset G$ such that $\gamma
K\subset G-A^G$. Since $ch_{G-A^G}=1$ on $\gamma K$, this condition is a
consequence of the condition for $D$ itself. It follows that $D_G$ is
admissible on $G-A^G$.

{\bf proof}: The result follows for $D_A$ provided for each $a\in A$ and for
some sufficiently small compact open subgroup $U_m\subset \subset A,$ and
each non-trivial character $\chi $ of $U_m$, we have 
\[
\begin{array}{c}
\int_{U_m}|D(k)|^{1/2}F_D(k)\chi (k)dk \\ 
=q^{-m}\int_{O_F}|D((1+\pi ^mu)a)|^{1/2}F_D((1+\pi ^mu)a)\chi (1+\pi ^mu)du
\\ 
=\sum_{n\geq m}d_nq^{-n}\int_{O_F^{\times }}F_D((1+\pi ^nu)a)\chi (1+\pi
^nu)du=0,
\end{array}
\]
where $d_n$ denotes the common value of $|D((1+\pi ^nu)a)|^{1/2}.$ (These
integrals exists since $|D(k)|^{1/2}F_D(k)$ is locally constant on the
regular set and bounded on $U_m$.) If $a\neq \pm 1$ then this follows from
the fact that $F_D$ is locally constant on the regular set.

This completes the proof. $\Box $

\subsection{Questions and an example}

{\sc Question}: Is is possible to classify the admissible representations $D$
in terms of the functions $F_D$? In other words, what conditions on $F_D$
guarantee that $D$ is admissible?

This question is related to the characterization of the algebraic dual of
the admissible Paley-Wiener space so would appear to be rather difficult.

{\sc Question}: Is is possible to classify which representable {\em tempered}
distributions are admissible? In other words, what conditions on the {\em %
tempered} functions $F_D$ guarantee that $D$ is admissible?

This question is related to the characterization of the algebraic dual of
the tempered Paley-Wiener space and so seems to be more tractable. In fact,
we will explore this question further in the case of $SL(2)$ below. As a
start, we prove the following

{\bf Lemma} {\em Let }$C\subset G_\sigma ^{\wedge }${\em \ }$\subset
G_t^{\wedge }${\em \ denote a compact Borel measurable set, where }$\sigma
\in M_c^{\wedge }${\em \ and }$M${\em \ is a Levi component of }$G,${\em \
and let }$d\lambda ${\em \ denote a bounded Borel measure on }$C${\em . The
distribution defined by} 
\[
D(f)=\int_C\Theta _\pi (f)d\lambda (\pi ),\;\;\;\;\;f\in C_c^\infty (G), 
\]
{\em exists and is admissible.}

{\bf proof}: By Harish-Chandra, each distribution $f\longmapsto \Theta _\pi
(f)$ is admissible. This implies that 
\[
(\Theta _\pi *\chi _\sigma ^0)(x)=0,\;{\rm on}\;U\;{\rm unless\;}G\;{\rm %
intertwines\;the\;trivial\;representation\;on\;}K\;{\rm with}\;\sigma , 
\]
for $\sigma ,\;K,\;U$ as above. Since $C$ is compact we may apply Fubini's
theorem to verify this admissibility condition for 
\[
F_D(x)=\int_C\Theta _\pi (x)d\lambda (\pi ).\; 
\]
This completes the proof. $\Box $

\subsection{Distributions on the tempered Paley-Wiener space, revisited}

Let $d\mu $ denote the Plancherel measure on the tempered dual $G_t^{\wedge
} $ and let $m(\pi )d\mu (\pi )$ denote a distribution on the tempered
Paley-Wiener space ${\cal C}(G)^{\wedge }$ such that

(1) $m(\pi )$ is supported on finitely many components of $G_t^{\wedge },$

(2) there is a continuous function $h$ on $G_t^{\wedge }$ such that on each
component $G_\sigma ^{\wedge },$ with $\sigma \in M_c^{\wedge }$ being a
cuspidal representation of some Levi $M$ of $G$, there is a function $%
h_\sigma $ on $X_{un}(M)$ such that $h(I_M(\omega \sigma ))$ is the
restriction of $h_\sigma (\omega )$ on the set 
\[
\{\omega \in X_{un}(M)\;|\;I_M(\omega \sigma )\in G_t^{\wedge }\}, 
\]
such that, as distributions on $\{\omega \in X_{un}(M)\;|\;I_M(\omega \sigma
)\in G_t^{\wedge }\},$ we have 
\[
m(I_M(\omega \sigma ))=\frac{\partial ^I}{\partial \omega ^I}h_\sigma
(\omega ), 
\]
where $I=(i_1,...,i_r)$ denotes a multi-index, $r$ being the real dimension
of $X_{un}(M)$, and $\frac{\partial ^I}{\partial \omega ^I}$ denotes partial
differentiation on the real manifold $X_{un}(M).$ Note that by (1) the
function $h$ in (2) is compactly supported. We call such a distribution a 
{\bf distribution of finite type} on $G_t^{\wedge }.$ A distribution
satisfying (2) but not (1) will be called a {\bf distribution of
quasi-finite type }on $G_t^{\wedge }.$ The maximum of the integers $%
|I|=i_1+...+i_r$, where $I$ runs over all multi-indices occurring in (2), is
called the {\bf order} of the distribution.

The theorem of section 6 implies the following

{\bf Theorem}: {\em If }$D${\em \ is a tempered distribution on }$G${\em \
then there is a distribution of quasi-finite type }$m(\pi )d\mu (\pi )${\em %
\ on the tempered Paley-Wiener space }${\cal C}(G)^{\wedge }${\em \ such that%
} 
\[
D(f)=\int_{G_t^{\wedge }}\Theta _\pi (f)m(\pi )d\mu (\pi ),\;\;\;\;\;f\in 
{\cal C}(G). 
\]

{\sc Question}: Is any such distribution is admissible?

By the previous lemma, the answer is yes for distributions of order $0$.

\subsection{Quasi-finite distributions and the Schwartz space}

We shall end this section with a proof, using a completely different
argument, of a similar statement for $H=GL(2,F)$. (We use $GL(2)$ rather
than $SL(2)$ since the presence of a non-compact center introduces a key
idea which will be expanded on and generalized in [J].) Let $Z$ denote the
center of $H$ for the remainder of this section.

{\bf Theorem}: {\em If }$D${\em \ is a tempered distribution on the space }$%
{\cal C}${\em $(H)$ then there is a distribution of quasi-finite type }$%
m(\pi )d\mu (\pi )${\em \ on tempered Paley-Wiener space }${\cal C}%
(H)^{\wedge }${\em \ such that} 
\[
D(f)=\int_{H_t^{\wedge }}\Theta _\pi (f)m(\pi )d\mu (\pi ),\;\;\;\;\;f\in 
{\cal C}(H). 
\]

{\bf proof} (Really, more of a sketch since some details are left out.): For
each character $\omega \in Z_u^{\wedge },$ consider the map ${\cal C}%
(H)\rightarrow {\cal C}_{*}(H,\omega )$ defined by $f\longmapsto f_\omega $,
where 
\[
f_\omega (x)=\int_Zf(xz)\omega (z)^{-1}dz, 
\]
and where 
\[
{\cal C}_{*}(H,\omega )=\{f\in {\cal C}_{*}(H)\;|\;f(xz)=\omega
(z)f(x),\;\;\;x\in H,\;z\in Z\}. 
\]
We have ${\cal C}(H)=\int_{\omega \in Z_u^{\wedge }}{\cal C}_{*}(H,\omega )$%
, so $D$ has a Fourier expansion 
\[
D(f)=\int_{\omega \in Z_u^{\wedge }}D_\omega (f_\omega )d\omega , 
\]
each $D_\omega $ being a distribution on ${\cal C}_{*}(H,\omega )$.
Following Kazhdan, let $W(\omega )$ denote the space of all invariant
tempered distributions on ${\cal C}_{*}(H,\omega )$ and, for each compact
open subgroup $K\subset \subset H$, let ${\cal C}_{*}(H,\omega )_K$ denote
the space of all linear combinations of translates of bi-$K$-invariant
functions in ${\cal C}_{*}(H//K,\omega \dot ),$ and let 
\[
{\cal C}_{*}(H,\omega )_K^{\perp }=\{f\in {\cal C}_{*}(H,\omega
)\;|\;\int_Kf(k_1kk_2)dk=0,\;\;k_1,k_2\in K\}. 
\]
By analogy with [K, section 4] (where it is assumed that the center is
compact), we have a direct sum decomposition ${\cal C}_{*}(H,\omega )={\cal C%
}_{*}(H,\omega )_K\bigoplus {\cal C}_{*}(H,\omega )_K$ and a corresponding
decomposition of the space of invariant distributions: $W(\omega )=W(\omega
)_K\bigoplus W(\omega )_K^{\perp }.$ We {\em claim} that $W(\omega
)=\bigcup_KW(\omega )_K$, but we shall not prove this here. (Argue by
contradiction ... ) By analogy with [K, section 4, Lemma 5], there is a
decomposition 
\[
D_\omega =\int_{G_{c,\omega }^{\wedge }}\Theta _\pi m_G(\pi )d\mu _\omega
(\pi )+D_\omega ^0, 
\]
where $D_\omega ^0$ is supported off the elliptic set, the distribution $m_G$
is a finite linear combination of Dirac delta functions, and the dual $%
G_{c,\omega }^{\wedge }$ is the set of equivalence classes of cuspidal
representations with unitary central character $\omega $. By the Weyl
integration formula, $D_\omega ^0$ must be supported on $A^H$ where $A$
denotes the diagonal subgroup of $H$. In the decomposition 
\[
D_\omega (f)=(D_\omega )_H(f)+(D_\omega )_A(F_f^A), 
\]
discussed in the Lemma of section 8.2 above (where $B$ here denoted the
upper triangular subgroup of $H$), we may identify $D_\omega ^0(f)$ with $%
(D_\omega )_A(F_f^A).$ The distribution $(D_\omega )_A$ is (tempered and)
admissible on $A$ by (analogy with) the Lemma in section 8.2. By the
Plancherel formula in the abelian case, $(D_\omega )_A$ may be represented
by a distribution of finite type on the tempered dual of $A$: 
\[
D_\omega ^0(f)=(D_\omega )_A(F_f^A)=\int_{A_t^{\wedge }}(F_f^A)^{\wedge
}(\chi )m_{A,\omega }(\chi )d\mu _A(\chi ). 
\]
On the other hand, by the character formulas for the principal series, $%
(F_f^A)^{\wedge }(\chi )=\Theta _{I_A(\chi )}(f),$ so 
\[
D_\omega ^0(f)=(D_\omega )_A(F_f^A)=\int_{A_t^{\wedge }}\Theta _{I_A(\chi
)}(f)m_A(\chi )d\mu _A(\chi ). 
\]
Therefore, 
\[
D=\int_{\omega \in Z_u^{\wedge }}D_\omega =\int_{\omega \in Z_u^{\wedge
}}\int_{G_{c,\omega }^{\wedge }}\Theta _\pi m_{G,\omega }(\pi )d\mu _\omega
(\pi )+\int_{\omega \in Z_u^{\wedge }}\int_{A_t^{\wedge }}\Theta _{I_A(\chi
)}m_{A,\omega }(\chi )d\mu _A(\chi ), 
\]
as desired.

This finishes the proof/sketch. $\Box $

\section{Supertempered distributions}

We close by discussing the supertempered distributions on $SL(2)$. The
definition of a supercuspidal distribution requires some preliminaries which
we present first.

\subsection{The constant term}

Let $P=MN$ denote the Levi decomposition and let $\overline{P}=M\overline{N}$
denote the Levi decomposition of the opposite parabolic. We shall consider
the maps ${\cal A}(G)\rightarrow {\cal A}(M),$defined by$\;f\longmapsto
f_P\;,\;_w{\cal A}(G)\rightarrow \;_w{\cal A}(M)$, defined by $f\longmapsto
_wf_P$. We say that a locally constant function $h$ on $G$ satisfies the 
{\bf weak inequality} if 
\[
|h(x)|<<\Xi (x)(1+\sigma (x))^r, 
\]
for some $r>0.$ Here $_w{\cal A}(G)$ denotes the space of matrix
coefficients satisfying the weak inequality.

For $t>0$, define $A^{+}(t)=\{a\in A\;|\;|\alpha (a)|\geq t\}$, where $%
\alpha \left( 
\begin{array}{cc}
a & 0 \\ 
0 & a^{-1}
\end{array}
\right) =a^2$, and let $A^{+}=A^{+}(1)$. Let $S\subset G$ be a set such that 
$S\cdot A^{+}(t_0)\subset S,$ for some $t_0>0$, and let $f:S\rightarrow V$
be a vector-valued function, where $V$ is some complex vector space. We say
that $f$ is {\bf negligible along} $P$, written $f\approx _P0$, if for every
compact subset $C\subset \subset G$ there is a $t\geq t_0$ such that $%
f(xa)=0 $ for all $x\in C\cap S,\;a\in A^{+}(t)$. Given two vector-valued
functions $f_i:S\rightarrow V,\;i=1,2$, we write $f_1\approx _Pf_2$ if $%
f_1-f_2\approx _P0.$ Let $V$ be a normed vector space with norm $||...||_V$.
For a fixed $x\in S$ and $L\in V$, we write 
\[
f(xa)\rightarrow _PL,\;a\rightarrow _P\infty ,\;\;\;\;\;or\;\;\;\;\;\;\lim
_{a\rightarrow _P\infty }f(xa)=L, 
\]
if for every $\varepsilon >0$ there is a $t_\varepsilon >0$ such that $%
||f(xa)-L||_V<\varepsilon $.

For each $f\in {\cal A}(G)$, define $f_{(B)}(a)=\delta _B(a)^{1/2}f(a)$, $%
a\in A$.

{\bf Theorem} (Harish-Chandra) {\em For each }$f\in {\cal A}(G)${\em , there
is a unique }$f_B\in {\cal A}(A)${\em , such that }$f_B\approx _Pf_{(B)}$%
{\em \ with }$S=A${\em .}

The function $f_B$ in this theorem is called the {\bf constant term of }$f$%
{\bf \ along }$B$. The construction of the constant term is in [S, section
2.6].

{\bf Theorem} (Harish-Chandra) {\em For each }$f\in _w{\cal A}(G)${\em ,
there is a unique }$_wf_B\in {\cal A}(A)${\em , such that }$\lim
_{a\rightarrow _P\infty }|_wf_B(aa^{\prime })-f_{(B)}(aa^{\prime })|=0${\em %
\ for all }$a^{\prime }\in A${\em .}

Following [S, section 1.10], we call a function $f$ (right) {\bf Hecke finite%
} if, for each compact open subgroup $K\subset \subset G$, the vector space
spanned by $f*g$, with $g\in C_c(G//K)$ (i.e., $g$ is compactly supported
and bi-$K$-invariant), is finite dimensional. We call an invariant
distribution $D$ on $G$ {\bf Hecke finite} if, for each $f\in C_c^\infty (G)$%
, the function $D_f(x):=D(r(x^{-1})f)$ is right Hecke finite as a function
of $x\in G$. Here $r$ denotes right translation: $(r(x)f)(g)=f(gx^{-1}).$
Examples of Hecke finite distributions are characters of admissible
representations of finite length (this is implicit in [S, section 1.11]). In
fact, a Hecke finite distribution is a virtual character ([J, section 3]),
hence automatically admissible.

{\bf Lemma}: {\em If }$D${\em \ is Hecke finite then }$D_f\in {\cal A}(G).$%
{\em \ If moreover }$D${\em \ is tempered then }$D_f\in \;_w{\cal A}(G)${\em %
.}

For each $f\in C_c^\infty (G)$, define $h_f\in C_c^\infty (M)$ by 
\[
h_f(m)=\delta _P(m)^{1/2}\int_N\int_{\overline{N}}f(\overline{n}mn)\delta _P(%
\overline{n})^{-1}d\overline{n}dn. 
\]

{\bf Proposition}: (1) {\em Let }$D${\em \ be a Hecke finite distribution.
For each parabolic }$P$, {\em there is an admissible distribution }$D_P${\em %
\ on }$M${\em \ such that } $(D_f)_P(1)=D_P(h_f)$, {\em \ for all }$f\in
C_c^\infty (G).$

(2) {\em Let }$D${\em \ be a Hecke finite tempered distribution. For each
parabolic }$P${\em , there is an admissible tempered distribution }$D_P${\em %
\ on }$M${\em \ such that }$_w(D_f)_P(1)=D_P(h_f),${\em \ for all }$f\in
C_c^\infty (G).$

We will prove this in [J]. The distribution $D_P$ on $M$ will be called the 
{\bf constant term of }$D${\bf \ along }$P$. This is not to be confused with
the $D_M$ introduced in section 7 above.

\subsection{Definition of a supertempered distribution}

Following Harish-Chandra in the real case, we say that a Hecke finite
admissible distribution $D$ is {\bf supertempered} if $D_P=0$ for all proper
parabolics $P$ of $G$. Examples of supertempered distributions are
characters of square-integrable representations (see [S, Corollary 4.4.7]).
Let $G_e$ denote the set of {\bf elliptic elements}: elements which are not
conjugate to an element in some proper parabolic subgroup of $G$. Other
examples of supertempered distributions arise from the lemma below:

{\bf Lemma}:{\em \ If }$D${\em \ is a Hecke finite distribution and supp}$%
(F_D)\subset G_e${\em \ then }$D${\em \ is supertempered and tempered.}

One proof of this, which is omitted, is to use the construction of the
constant term to verify directly that $(D_f)_P(1)=0$, for all proper
parabolics $P$ (see [J, section 2]). The temperedness is an immediate
consequence of a result of Kazhdan [K, section 5, Lemma 3].

{\sc Question} If $D$ is supertempered and ${\em F_D|_{G_e}=0}$ then does it
necessarily follow that $D=0${\em ?}

{\sc Question}: Does supertempered imply tempered?

{\bf Lemma}: {\em If }$D${\em \ is a Hecke finite distribution on }$G${\em \
then }$D${\em \ is a virtual character.}

This is not especially hard to prove (see [J, section 3]).

{\bf Corollary}: {\em If }$D${\em \ is a supertempered distribution then }$D$%
{\em \ is a virtual character.}

\subsection{An example}

Let $\chi =sign_\tau ,$ for some$\;\tau =\varepsilon ,\pi ,\pi \varepsilon $%
, and let $\pi =T_\chi ^{+}$,\ $\pi ^{\prime }=$ $T_\chi ^{-}.$

{\bf Proposition}: {\em The distribution } 
\[
D=\Theta _\pi -\Theta _{\pi ^{\prime }} 
\]
{\em is supertempered. }

{\bf proof} The character computation of section 4.2.2 above shows that $D$
is supported in the elliptic set, so by the above lemma it must be
supertempered. It is instructive to give a second proof directly from the
definition.

Consider the matrix coefficients 
\[
f_1(x)=\int_g\Theta _\pi (g)h(gx^{-1})dg,\;\;f_2(x)=\int_g\Theta _{\pi
^{\prime }}(g)h(gx^{-1})dg,\;\;\;h\in C_c^\infty (G), 
\]
and let 
\[
\begin{array}{c}
f_{1,(B)}(a)=\delta _B(a)^{1/2}f_1(a)=\delta _B(a)^{1/2}\int_G\Theta _\pi
(g)h(ga^{-1})dg, \\ 
f_{2,(B)}(a)=\delta _B(a)^{1/2}f_2(a)=\delta _B(a)^{1/2}\int_G\Theta _{\pi
^{\prime }}(g)h(ga^{-1})dg.
\end{array}
\]
The character formulas above imply that 
\[
\begin{array}{c}
f_{1,(B)}(a)-f_{2,(B)}(a)=\delta _B(a)^{1/2}\int_G(\Theta _\pi (g)-\Theta
_{\pi ^{\prime }}(g))h(ga^{-1})dg \\ 
=\delta _B(a)^{1/2}\sum_{T\neq A}\frac 12\int_T|D(t)|^{1/2}(\Theta _\pi
(t)-\Theta _{\pi ^{\prime }}(t))h(ta^{-1})dt.
\end{array}
\]
In fact, the sum over the non-split Cartans $T$ runs only over one or two
Cartans. To finish the proof of this proposition, we use the following

{\bf Lemma}: {\em If }$h${\em \ is supported on }$C\subset \subset G${\em ,
there exists an }$r=r_C>0${\em \ such that, for all Cartans }$T${\em \ in
the above sum and all }$a\in A^{+}(r)${\em , we have }$C\cap aT=\emptyset .$

{\bf proof}: To see this, note that there is an $N>0$ such that $\sigma
(x)<N $ for all
\[
x\in C\cup \{T\;|\;T\neq A\;{\rm occurs\;in\;the\;sum\  above}\}.
\]
This is because $C$ is compact and each $T$ is compact. On the other
hand, for each Cartan $T\neq A$, from the explicit description of $T$ as a
matrix group ([SS2]) and the definition of $\sigma (x)$, it follows that
there is an $r>0$ (depending on $N,\;T$) such that $\sigma (ta)>N$ for all
$a\in A^{+}(r),\;t\in T$. From this the lemma follows. $\Box (lemma)$

From the lemma and the definition of negligible, it follows that
$f_{1,(B)}(a)-f_{2,(B)}(a)\approx _B0$ along $A$. This implies that $D_B=0$.
$\Box (proposition)$

Combining a special case of a result in [H] along with the above corollary
in section 9.2, we find

{\bf Theorem}: (Herb) A tempered, supertempered distribution of $SL(2,F)$ is
a finite linear combination of discrete series characters and the
distributions in the previous proposition.

\section{References}



\begin{thebibliography}{99}


\bibitem[A]{A}
  %[A]
  M. Assem, ``The Fourier transform and some character formulae for p-adic 
$SL(\ell )$, $\ell $ a prime'', preprint

\bibitem[Ba]{Ba}
  %[Ba]
  G. Bachman, {\bf Elements of abstract harmonic analysis},
  Academic Press, New York, 1964

\bibitem[B]{B}
  %[B]
  J. Bernstein, ``All reductive p-adic groups are tame'', Funct. Anal.
Appl. 8(1974)91-93

\bibitem[BDK]{BDK}
  %[BDK]
  J. Bernstein, P. Deligne, D. Kazhdan, ``Trace Paley-Wiener theorem for
reductive p-adic groups'', J. d'Anal. Math. 47(1986)180-192

\bibitem[Ca]{Ca}
  %[Ca]
  W. Casselman, ``Introduction to the theory of admissible
representations on reductive p-adic groups'', preprint

\bibitem[C]{C}
  %[C]
  L. Clozel, ``Invariant harmonic analysis on the Schwartz space of a
reductive p-adic group'', in {\bf Harmonic analysis on reductive groups}, ed. W. Barker, P.
Sally, Birkh\"auser, Boston, 1991

\bibitem[D]{D}
  %[D]
  J. Dixmier, {\bf $C^{*}$-algebras}, North-Holland, Amsterdam, 1977

\bibitem[F]{F}
  %[F]
  J. Fell, ``$C^*$-algebras with smooth dual'', Ill. J. Math.
4(1960)221-230

\bibitem[Fr]{Fr}
  %[Fr]
  S. Franklin, ``The reducible principal series of SL(2) over a p-adic
field'', Thesis, Univ. Chicago, 1971

\bibitem[GG]{GG}
  %[GG]
  I. M. Gelfand, M. I. Graev, ``Representations of a group of the second
order with elements from a locally compact field'', Russian Math. Surveys,
18(1963)29-100

\bibitem[H]{H}
  %[H]
  R. Herb, ``Virtual supertempered characters'', preprint

\bibitem[H1]{H1}
  %[H1]
  L. H\"ormander, {\bf An introduction to complex analysis in several
    complex variables}, 2nd edition, North-Holland,
Amsterdam, 1973

\bibitem[H2]{H2}
  %[H2]
  -----, {\bf The analysis of linear partial differential operators, I},
  Springer-Verlag, New York, 1983

\bibitem[J]{J}
  %[J]
  D. Joyner, ``Some results concerning invariant distributions on a
reductive group over p-adic field'', preprint, 1994

\bibitem[K]{K}
  %[K]
  D. Kazhdan, ``Cuspidal geometry of p-adic groups'', J. d'Anal. Math.
47(1986)1-36

\bibitem[M]{M}
  %[M]
  G. Mackey, ``Borel structure in groups and their duals'', Trans. AMS
85(1957)134-165

\bibitem[S]{S}
  %[S]
  A. Silberger, {\bf Introduction to harmonic analysis on reductive
    p-adic groups}, Princeton Univ. Press, Princeton, 1979

\bibitem[Sch]{Sch}
  %[Sch]
  L. Schwartz, {\bf Theorie des distributions, vol. I},
  Hermann, Paris, 1950

\bibitem[SS1]{SS1}
  %[SS1]
  P. Sally, J. Shalika, ``Characters of the discrete series of
representations of SL(2) over a p-adic field'', Proc. Acad. Sci. USA
61(1968)1231-1237

\bibitem[SS2]{SS2}
  %[SS2]
  ------, ``The Fourier transform of orbital integrals on SL(2) over a
  p-adic field'', in {\bf Lie group representations, II},
  ed., R. Herb et al, Lecture Notes in Math. vol 1041,
Springer-Verlag, New York, 1984

\bibitem[T]{T}
  %[T]
  M. Tadi\'c, ``The topology of the dual space of a reductive group over a
local field'', Glasnik Matem., 18(1983)259-279

\bibitem[Ta]{Ta}
  %[Ta]
  M. Taibleson, {\bf Fourier analysis on local fields},
  Princeton Univ. Press, Princeton, 1975
 
\end{thebibliography}
D. Joyner, Math Dept, USNA, Annapolis, MD 

\end{document}
